\chapter{Verification, dead ends, numerics and open problems}\label{ch:15}

\providecommand{\chXVsq}{\mathbin{\ltimes}}

This chapter documents the release of the library named in the front matter; every Lean quotation is checked by \texttt{scripts/textbook/sync.py} to occur verbatim (up to whitespace and
comments) in the source files of that release. Its body was first written against the commit ``Construct the
Moonshine module $V^\natural$ (FLM) unconditionally'' (7 October 2026, then the head of the main branch; ``the Moonshine
commit'' below); records tied to that snapshot, and other checks not re-run for this edition, are marked as such, and the
update of October 2026 is described below. The
preceding fourteen chapters presented the mathematics of the LeanVOA library and, as their
culmination, the construction of the Moonshine module $\Vnat$. Each of them followed the same
discipline: a mathematical statement, the Lean declaration copied verbatim from the repository, a
fidelity check explaining why the one is the other, an argument on paper, and an honest status
report. This last chapter turns the discipline on the library as a whole. It asks four questions
that a referee-editor must be able to answer before signing off on the book.

\begin{enumerate}[label=(\arabic*)]
\item \emph{What exactly has been checked, by what, and how would we know if it had not?} The
  answer is a verification protocol made of a compiler run, two scripts and a manual audit, each
  of which certifies something specific and nothing more (Section~\ref{ch15:sec:protocol}).
\item \emph{Which statements were wrong, vacuous or circular at some point, and how was that
  discovered?} The project kept a register of dead ends, traps and corrected mistakes, library-wide
  and for the Moonshine construction in particular. Section~\ref{ch15:sec:deadends} consolidates
  it, each entry with its precise statement, the certificate of failure and the lesson.
\item \emph{What was computed numerically, and what is the status of those computations?}
  Section~\ref{ch15:sec:numerics} describes every numerical experiment of the Moonshine
  construction, with its method, tolerance and controls, and says plainly that none of it is part
  of any proof.
\item \emph{Why was the construction done this way, can it be reproduced, and what is still
  open?} Sections~\ref{ch15:sec:design}, \ref{ch15:sec:repro} and~\ref{ch15:sec:open}.
\end{enumerate}

Throughout, four registers of claims are kept strictly apart, as the project's working agreements
demand:
\begin{itemize}
\item \textbf{proved unconditionally} in Lean: the declaration has no hypothesis that is not
  itself established, and \texttt{\#print axioms} lists only Lean's three standard axioms;
\item \textbf{proved conditionally} in Lean: the declaration is accepted by the kernel, but its
  statement carries a hypothesis (a \lean{Prop}-valued premise, or a structure that has to be
  supplied) which is not proved anywhere; the hypothesis is then named;
\item \textbf{checked numerically}: a computer calculation outside Lean, in floating point or in
  exact arithmetic, agrees with the claim on a finite sample; this is evidence, never proof;
\item \textbf{argued on paper only}, or \textbf{open}.
\end{itemize}
Earlier editions had a fifth register, \textbf{proved using the project axiom}
\lean{VOA.pbw}; since 19 October 2026 that axiom is a theorem, and the register is empty
(Section~\ref{ch15:sec:trust}; Chapter~\ref{ch:L02}, Section~\ref{l02:sec:history}).
The Moonshine construction never used it.

\paragraph{The update of October 2026.} The body of this chapter was written against the Moonshine commit. In October 2026
the book was brought up to date with the library on the branch \lean{deligne-lean}: five chapters were added
(Chapters~\ref{ch:L23}--\ref{ch:L27}: affine vertex algebras for every simple Lie algebra and the four chapters on the
Deligne series; the chapter on affine $\mathfrak{sl}_3$ is in preparation; the two chapters on graded unitarity through the
BRST coset, Chapters~\ref{ch:L28} and~\ref{ch:L29}, were added for the v1.0 release), Chapters~\ref{ch:L11}, \ref{ch:L19}
and~\ref{ch:L20} were extended, and every chapter that described \lean{VOA.pbw} as an axiom was corrected. The present
chapter was extended, not rewritten, to cover that work. The additions are: the book's own verification record and the
trust base at the manifest commit (Section~\ref{ch15:sec:bookrecord}); the independent release audits of
27--28~October~2026 (Section~\ref{ch15:sec:releaseaudits}); the dead ends recorded in the new and extended chapters
(Section~\ref{ch15:sec:deadends-oct}); and the open problems that those chapters name
(Section~\ref{ch15:sec:open-later}). Following the book's scope decision, work in progress in the library is not
described; where a route is under development, we say so and nothing more.

\begin{remark}[A note on the physics]\label{ch15:rem:physics}
A physicist reading this chapter may wonder why so much effort goes into checking that a
statement means what it says. In a conformal field theory one rarely writes down the axioms of
the operator algebra completely; one computes correlation functions and trusts that locality,
associativity and modular invariance hold because the theory ``exists''. A formalization cannot
work that way. Every axiom is a line of code, and a mistyped quantifier produces a theory that
either does not exist (no nonzero object satisfies the axioms) or exists for trivial reasons (the
zero theory, or a degenerate extension). Both happened in this project. The checks of this chapter
are the mathematical counterpart of a physicist's sanity checks: is the partition function right,
is the vacuum unique, does the operator product expansion close, and is the theory non-trivial.
\end{remark}

% ---------------------------------------------------------------------------------------------
\section{The verification protocol}\label{ch15:sec:protocol}

\subsection{What a Lean proof certifies, and what it does not}\label{ch15:sec:kernel}

A Lean file is a sequence of declarations: definitions (\lean{def}, \lean{structure},
\lean{abbrev}) and theorems (\lean{theorem}). A theorem declaration has the shape
\begin{plaincode}
theorem name (hypotheses) : statement := proof
\end{plaincode}
When the file is compiled, Lean's \emph{kernel}, a small and carefully audited type checker,
verifies that \lean{proof} is a term of type \lean{statement} in the context of the hypotheses.
This is the only thing a successful compilation certifies, and it is a great deal: every step of
every argument has been checked down to the rules of dependent type theory. It is not everything a
reader wants, for three reasons.

\begin{enumerate}[label=(\alph*)]
\item \emph{Definitions are trusted, not verified.} If \lean{statement} mentions a notion such as
  \lean{VertexOperatorAlgebra}, \lean{IsSimple} or \lean{TwistedModule}, the kernel checks the
  proof against the definition as written. It cannot know whether the definition captures the
  intended mathematical notion. A definition with a mistyped quantifier is a perfectly good
  definition of something else. This is why the book prints every definition verbatim and checks
  its fidelity; it is also the most common source of the dead ends of
  Section~\ref{ch15:sec:deadends}.
\item \emph{Hypotheses are assumed, not proved.} A theorem with a hypothesis
  \lean{(h : P)} proves ``$P$ implies the statement''. If $P$ is false, or unsatisfiable, the
  theorem is true and worthless.
\item \emph{Axioms are assumed.} A declaration introduced with \lean{axiom} is accepted without
  proof, and every theorem whose proof term mentions it, directly or through a chain of lemmas,
  depends on it.
\end{enumerate}

The verification protocol of the project addresses each of these. Point~(c) is handled
mechanically, by \texttt{\#print axioms} and the axiom audit; point~(b) by a rule of language
(``unconditional'' versus ``conditional'') and a manual trace of hypotheses; point~(a) only by human
reading, which is what the fidelity checks of this book are for.

\subsection{The trust base: three standard axioms (and the former project axiom)}\label{ch15:sec:trust}

Lean's logic has three axioms that essentially all of Mathlib uses:
\begin{itemize}
\item \lean{propext} (\emph{propositional extensionality}): two propositions that imply each other
  are equal;
\item \lean{Quot.sound}: the quotient of a type by a relation identifies related elements; this
  is what makes quotient constructions (such as quotient modules, or the Zhu algebra
  $A(V)=V/O(V)$) behave as in ordinary mathematics; it also implies function extensionality;
\item \lean{Classical.choice}: from a proof that a type is nonempty one may extract an element;
  this gives the law of excluded middle and all of classical reasoning, including the choice of
  bases of vector spaces.
\end{itemize}
These three are the background logic of ordinary mathematics (Zermelo--Fraenkel set theory with
choice proves corresponding statements), and a theorem whose axiom list is a subset of
$\{\lean{propext},\lean{Classical.choice},\lean{Quot.sound}\}$ is called
\emph{\texttt{\#print axioms}-clean}.

The library adds no axiom of its own. Until 19 October 2026 it added exactly one, in
\leanpath{VOA/Foundations/AffineLie.lean}; the same declaration is now a theorem:
\begin{leancode}
structure PBWPrinciple : Prop where
  /-- PBW nonvanishing for every affine vacuum module. -/
  affineVacuum : ∀ {𝔤 : Type u} [LieRing 𝔤] [LieAlgebra ℂ 𝔤]
      (κ : LinearMap.BilinForm ℂ 𝔤) (hκ : κ.IsInvariant) (k : ℂ),
    PBW κ hκ k
  /-- The canonical bottom embeds in every module induced from a Lie subalgebra. -/
  inducedBottom : ∀ {g U : Type u} [LieRing g] [LieAlgebra ℂ g]
      [AddCommGroup U] [Module ℂ U] (p : LieSubalgebra ℂ g)
      (rho : p →ₗ⁅ℂ⁆ Module.End ℂ U),
    Function.Injective (PBWInduction.bottom p rho)

theorem pbw : PBWPrinciple :=
  ⟨fun κ hκ k => AffineLie.vacuumVector_ne_zero_of_induction κ hκ k,
    fun p rho => PBWTheorem.bottom_injective p rho⟩
\end{leancode}
where \lean{PBW κ hκ k} is defined, in the same file, as
\lean{AffineLie.vacuumVector κ hκ k ≠ 0}.

\begin{dictionary}
\textbf{What the principle says.} The Poincar\'e--Birkhoff--Witt theorem says that the universal
enveloping algebra $U(\mathfrak g)$ of a Lie algebra has a basis of ordered monomials in a basis of
$\mathfrak g$; consequently $U(\mathfrak g)$ is free as a right module over $U(\mathfrak p)$ for any
subalgebra $\mathfrak p$. The principle bundles two consequences of PBW:
\begin{itemize}
\item \lean{affineVacuum}: for every Lie algebra $\mathfrak g$ over $\C$ with an invariant
  bilinear form $\kappa$ and every level $k\in\C$, the vacuum vector of the affine vacuum module
  $V^k(\mathfrak g)=U(\hat{\mathfrak g})\otimes_{U(\mathfrak g[t]\oplus\C K)}\C_k$ is nonzero.
  (Without PBW one cannot exclude that the induced module collapses to $0$.)
\item \lean{inducedBottom}: for a Lie subalgebra $\mathfrak p\subseteq\mathfrak g$ and a
  representation $\rho$ of $\mathfrak p$ on $U$, the canonical map
  $U\to U(\mathfrak g)\otimes_{U(\mathfrak p)}U$, $u\mapsto 1\otimes u$, is injective.
\end{itemize}
Both are classical theorems; Mathlib (v4.30.0) does not provide PBW, and the library now proves
the needed relative form itself (Chapter~\ref{ch:L02}, Section~\ref{l02:sec:pbwproof}). The
quantification over all universes \lean{Type u} is the
reason the statement is a \lean{structure} of two fields rather than two separate declarations.
\textbf{Could it be mis-stated?} Since it is proved, it cannot be inconsistent; the remaining
risk is a mis-encoding of \lean{PBWInduction.bottom} or \lean{vacuumVector} that would make it
weaker than intended, which is why
those definitions are printed and checked in Section~\ref{l02:sec:axiom}.
\end{dictionary}

\begin{physics}
PBW is the algebraic statement that a Verma-type module built by acting with creation operators on
a highest-weight state is ``as large as it looks'': the states
$J^{a_1}_{-n_1}\cdots J^{a_r}_{-n_r}|0\rangle$ with ordered modes are linearly independent before any
null vectors are quotiented out. Every physicist uses this silently when counting the states of a
current algebra (the Wess--Zumino--Witten vacuum module at level $k$) or when writing the character
$\prod_n(1-q^n)^{-\dim\mathfrak g}$ of the universal affine algebra. \lean{VOA.pbw} (with \lean{pbwBasis_holds}) proves exactly
that counting, nothing more.
\end{physics}

\paragraph{Where \texorpdfstring{\texttt{VOA.pbw}}{VOA.pbw} is used, and the twin pattern (history).}
While \lean{VOA.pbw} was an axiom, the library followed a strict pattern (Chapter~\ref{ch:L02},
Section~\ref{l02:sec:history}): wherever a result needed PBW, it was proved twice. One version carries the needed consequence as an explicit hypothesis (a
datum such as \lean{Pr : BottomProjection ...}) and is \texttt{\#print axioms}-clean; a companion
wrapper, whose name ends in \lean{_voaPBW} or \lean{_of_voaPBW}, supplies that datum from
\lean{VOA.pbw}. Since \lean{VOA.pbw} became a theorem the datum is supplied by a theorem, so each
\lean{_voaPBW} wrapper is the form of its twin with the PBW datum discharged. The two most
important instances are Zhu's bijection and Zhu's modularity theorem. From \leanpath{VOA/Modules/ZhuVermaAnArchLUBijection.lean} and
\leanpath{VOA/Modules/ZhuVermaAnArchLUBijectionPBW.lean}:
\begin{leancode}
theorem zhuBijective (Pr : BottomProjection (parabolicRepOfAMod W (regularAMod W))) :
    Function.Bijective (forwardMapIrr W) :=

theorem zhuBijective_of_voaPBW (W : VertexOperatorAlgebra 𝒱) :
    Function.Bijective (forwardMapIrr W) :=
  zhuBijective
    (voaPBWBottomProjection (parabolicRepOfAMod W (regularAMod W)))
\end{leancode}
and from \leanpath{VOA/Modularity/GNProp8Closure.lean} and \leanpath{VOA/Modularity/GNProp8PBW.lean}:
\begin{leancode}
theorem zhuModularity_pinned_of_C2Cofinite [Finite (Φ W)]
    (Mfam : ι → IrrModule W) (h : ι → ℂ)
    (hC2 : C2Cofinite W)
    (hss : letI := instRingA W; IsSemisimpleRing (A W))
    (Pr : BottomProjection (parabolicRepOfAMod W (regularAMod W)))
    (hcard : Fintype.card ι = Nat.card (Φ W))
    (hcw : ∀ i, HasConformalWeight (Mfam i).isMod (h i))
    (hsep : ∀ j k : ι, j ≠ k → h j = h k →
      IsEmpty (AModEquiv W (lowestAMod (hcw j).2) (lowestAMod (hcw k).2))) :
    ∃ D : ModularData W ι, D.Mfam = Mfam ∧ D.h = h ∧ ZhuModularityAnalytic D :=

theorem zhuModularity_pinned_of_C2Cofinite_voaPBW [Finite (Φ W)]
    (Mfam : ι → IrrModule W) (h : ι → ℂ)
    (hC2 : C2Cofinite W)
    (hss : letI := instRingA W; IsSemisimpleRing (A W))
    (hcard : Fintype.card ι = Nat.card (Φ W))
    (hcw : ∀ i, HasConformalWeight (Mfam i).isMod (h i))
    (hsep : ∀ j k : ι, j ≠ k → h j = h k →
      IsEmpty (AModEquiv W (lowestAMod (hcw j).2) (lowestAMod (hcw k).2))) :
    ∃ D : ModularData W ι, D.Mfam = Mfam ∧ D.h = h ∧ ZhuModularityAnalytic D :=
\end{leancode}

\begin{dictionary}
\textbf{The pattern, decoded.} The two modularity statements are identical except that the first
has the extra binder \lean{(Pr : BottomProjection ...)}. A \lean{BottomProjection} for the
parabolic representation attached to the regular $A(W)$-module is a linear projection from the
induced module onto its bottom piece, compatible with the structure; it exists if and only if the
bottom embeds (the \lean{inducedBottom} consequence of PBW). So the first theorem says ``if such a
projection exists, Zhu's modularity holds'', the second says ``Zhu's modularity holds'', and the
second is the first applied to the projection \lean{voaPBWBottomProjection}, which is derived
from \lean{VOA.pbw}. For the referee this means: the content of Zhu's theorem is entirely in the
\texttt{\#print axioms}-clean version, and \lean{VOA.pbw} enters only at one clearly marked seam. The
remaining hypotheses (\lean{hC2}: $C_2$-cofiniteness; \lean{hss}:
semisimplicity of $A(W)$, which follows from (admissible) rationality by
\lean{VOA.Huang.ZhuSemisimple.isSemisimpleRing_A_of_admissiblyRational}, \leanpath{VOA/Huang/ZhuSemisimple.lean}; \lean{hcard}, \lean{hcw},
\lean{hsep}: the family \lean{Mfam} is a complete irredundant list of irreducible modules with
conformal weights \lean{h}) are the hypotheses of Zhu's theorem itself; Chapter~\ref{ch:L15}
explains each. (Earlier versions also carried \lean{hzero} and \lean{hsimp}; both were redundant and have been removed.)
\end{dictionary}

\begin{status}
Re-run for this chapter (\lean{lake env lean -DmaxSynthPendingDepth=3} on a scratch file
importing \lean{VOA}, at the Moonshine commit; not re-run for this edition):
\lean{zhuModularity_pinned_of_C2Cofinite} and \lean{zhuBijective}: \lean{[propext,
Classical.choice, Quot.sound]}. \lean{zhuModularity_pinned_of_C2Cofinite_voaPBW},
\lean{zhuBijective_of_voaPBW}, \lean{AffineLie.vacuumVector_ne_zero} and
\lean{VHat.DLM3.voaPBWBottomInjective}: \lean{[propext, Classical.choice, Quot.sound, VOA.pbw]} (at
that commit, when \lean{VOA.pbw} was an axiom). Re-probed at a later commit, after \lean{VOA.pbw} became a theorem, all four
reported \lean{[propext, Classical.choice, Quot.sound]} (not re-run for this edition).
The same probe reports \lean{[propext, Classical.choice, Quot.sound]} for
\lean{heisenbergVOA}, \lean{VirModule.virasoroVOA}, \lean{AffineVOA.affineSl2VOA_closed},
\lean{LatGen.latticeVOAGen}, \lean{levelOne_coe_mem_E46}, \lean{LQ43.L43_graded_rigidity},
\lean{finiteDimensional_A_of_C2Cofinite}, \lean{onePointSubexp_of_c2},
\lean{strictSpanning_of_C2Cofinite}, \lean{simpleAModsRealized_of_C2Cofinite},
\lean{VertexOperatorAlgebra.borcherds}, \lean{Moonshine.LatticeTwisted.leechTwistedModule},
\lean{Moonshine.sectorSeries_eq_JShift} and \lean{Moonshine.moonshine}.
In particular, the affine $\mathfrak{sl}_2$ vertex operator algebra and the whole Moonshine chain
do not depend on \lean{VOA.pbw}.
\end{status}

Many files mention the string \lean{pbw} (a word search finds it in several hundred files, mostly
in documentation and in names such as \lean{PBWFreeness}); this says nothing about dependence.
Dependence is a property of proof terms and is decided only by \texttt{\#print axioms}.

\subsection{What \texorpdfstring{\lean{lake build}}{lake build} certifies: aggregate versus isolated}\label{ch15:sec:build}

\lean{lake} is Lean's build tool. The library is one Lean package, \lean{VOA}, whose root module
\leanpath{VOA.lean} (the \emph{umbrella}) imports the category umbrellas
\leanpath{VOA/Foundations.lean}, \leanpath{VOA/Core.lean}, \ldots, \leanpath{VOA/Moonshine.lean},
which in turn import every production file. Running \lean{lake build} with no argument builds the
default target \lean{VOA}, i.e.\ compiles the umbrella and, transitively, everything it imports,
together with the needed part of Mathlib. The output ends with a line such as
\begin{plaincode}
Build completed successfully (10343 jobs).
\end{plaincode}
where a \emph{job} is one unit of work (compiling one module, or one of its artefacts). The job
count is a useful fingerprint: if two runs on what should be the same tree give different counts,
the trees differ.

Two kinds of build must be distinguished.
\begin{itemize}
\item An \textbf{isolated build}, \lean{lake build VOA.Moonshine.TTBlock}, compiles one module and
  its imports. It certifies that this file, in the presence of exactly the files it imports, is
  correct.
\item The \textbf{aggregate build}, \lean{lake build} (equivalently \lean{lake build VOA}),
  compiles the umbrella. It certifies that \emph{all} production files compile \emph{together}.
\end{itemize}
The difference matters. A file can compile in isolation and still break the aggregate: two files
may declare the same name in the same namespace (the second import then fails), a file may add an
instance or a \lean{simp} lemma that changes the behaviour of proofs in another file that imports
both, or a file may simply not be imported by the umbrella at all, in which case its theorems are
not part of the library that the umbrella describes. During the Moonshine work many files were first
built in isolation by background agents (for example the route notes record
``\lean{lake build VOA.Moonshine.TTBlock} is green (4243 jobs)'' and ``Build completed successfully
(9039 jobs)'' for three \lean{SigmaExist} files) and only later wired into
\leanpath{VOA/Moonshine.lean}. The project rule is therefore:

\begin{quote}
\emph{A result counts only when the aggregate \lean{lake build} is green.}
\end{quote}

A second, subtler point was learned the hard way earlier in the library's history: a green build is a property of a \emph{tree at an instant}, not of a session.
Twice, a recorded ``green at $N$ jobs'' failed to reproduce because a later, unrelated file had been
added to the working copy, and once because a second session was editing files while the build ran.
Hence the operational rules: re-run the aggregate build immediately before recording a result,
record the job count actually seen, never run two builds concurrently in one working tree, and
prefer separate \lean{git worktree}s for concurrent work.

\subsection{\texorpdfstring{\texttt{scripts/check\_repository.py}}{The repository check}}\label{ch15:sec:checkrepo}

The aggregate build cannot see files that nobody imports, and it cannot forbid tokens that Lean
accepts. Both gaps are closed by a Python script, \lean{scripts/check_repository.py}, which is also
run in continuous integration before the Lean build. It does three things.

\begin{enumerate}
\item \textbf{Lexing.} It tokenizes every production file \lean{VOA/**/*.lean} and the umbrella,
  correctly skipping line comments \lean{--}, nested block comments \lean{/- ... -/} and string
  literals, so that the word ``sorry'' in a docstring is not a false positive.
\item \textbf{Forbidden tokens and the sanctioned axiom.} Its policy is the constant
\begin{plaincode}
ALLOWED_AXIOM_FILE = SOURCE_ROOT / "Foundations" / "AffineLie.lean"
FORBIDDEN_TOKENS = frozenset(
    {"sorry", "admit", "native_decide", "stop", "implemented_by"}
)
\end{plaincode}
  Any code occurrence of these five tokens is an error. They are, respectively: the placeholder
  for a missing proof; its tactic synonym; a tactic that evaluates a decidable proposition by
  compiling it to machine code (thereby trusting the compiler, not only the kernel); the tactic
  that abandons the rest of a proof; and an attribute that replaces a definition's executable code
  by other code. Since \lean{VOA.pbw} became a theorem every occurrence of the keyword \lean{axiom} is
  an error (the constant \lean{ALLOWED_AXIOM_FILE} survives in the script but no axiom is
  accepted); the script still prints a count of ``sanctioned axiom declarations'', now 0. (At
  the Moonshine commit the one exception was \lean{axiom pbw} in \leanpath{VOA/Foundations/AffineLie.lean}.)
\item \textbf{Reachability.} Starting from \leanpath{VOA.lean} it follows every local import
  (\lean{import VOA....}) and reports (i) imports naming a non-existent file and (ii) production
  files not reachable from the umbrella.
\end{enumerate}
Its success line, as run for the first edition of this chapter on the tree of the Moonshine commit, was
\begin{plaincode}
Repository check passed: 1867 library files scanned; 1867 reachable through 4676 local import edges; 1 sanctioned axiom declaration(s).
\end{plaincode}
So, at that commit, every one of the 1867 files then in the library was part of the aggregate build, none contained a
forbidden token, and there was exactly one axiom. At a later commit, after \lean{VOA.pbw} became a theorem, the script reported ``7107
library files scanned; 7105 reachable through 17543 local import edges; 0 sanctioned axiom
declaration(s)'', with verdict ``failed'' only because two draft files in \leanpath{VOA/Huang/} are
unreachable. The reachability check was the one that fired most often during the
Moonshine work. Agents wrote new files that compiled in isolation; until a maintainer added the
import line, the check listed them as unreachable, and on one occasion it caught an in-flight,
unwired draft (\leanpath{VOA/Moonshine/HeisenbergTwistedModuleFieldComm.lean}) that still contained
\lean{sorry}. That file was not committed in that state; it was later completed and wired.

\subsection{\texorpdfstring{\texttt{scripts/audit/run\_axiom\_audits.py}}{The axiom audit} and its blind spot}\label{ch15:sec:audit}

The third mechanical check, \lean{scripts/audit/run_axiom_audits.py}, enforces \emph{exact} axiom
sets for a maintained list of declarations. It runs 20 Lean ``probe'' files under
\leanpath{scripts/audit/} (for example \lean{RigidityProgressAxiomAudit.lean},
\lean{Phase5AxiomAudit.lean}, \lean{SuperAxiomAudit.lean}), each of which consists of
\texttt{\#print axioms} commands; it parses Lean's answers and compares each with the expected set:
by default $\{\lean{propext},\lean{Classical.choice},\lean{Quot.sound}\}$; for the single
declaration \lean{VOA.zhuModularity_pinned_of_C2Cofinite_voaPBW} the default plus \lean{pbw} (while
\lean{VOA.pbw} was an axiom; since it became a theorem the expected set \lean{PBW_AXIOMS} equals the default,
so any reappearance of \lean{pbw} is a failure); and for
a short explicit list of elementary lemmas a \emph{smaller} set (some use only \lean{propext}, or
no axiom at all). It fails if a declaration is missing, reported twice, or has a different set,
larger \emph{or smaller}. Insisting on exactness in both directions is deliberate: a set that
shrinks unexpectedly usually means that a declaration was replaced by a different one with the same
name. The success line of the final run was
\begin{plaincode}
Axiom audit enforcement passed: 4876 reports across 20 probes.
\end{plaincode}

\begin{warning}
\textbf{The blind spot.} The 20 probes were written before the Moonshine construction began, and
none of them mentions a declaration under \lean{VOA.Moonshine} or \lean{VOA.LatticesAndCodes}. The
report count 4876 is \emph{identical} before and after the construction (it is the count recorded
in the project's build records at each of the last four checkpoints of the Moonshine work). A green audit therefore says nothing
whatsoever about $\Vnat$. The headline Moonshine theorems were audited individually instead: at each
checkpoint a scratch file importing \lean{VOA} and containing \texttt{\#print axioms} for the
milestone declarations was compiled with \lean{-DmaxSynthPendingDepth=3}
(Section~\ref{ch15:sec:trap}), and the result was recorded in the project's internal records and in the route
notes. Adding a Moonshine probe to the maintained list is an open housekeeping task
(Section~\ref{ch15:sec:open}).
\end{warning}

\subsection{\texorpdfstring{\texttt{\#print axioms}}{print axioms} and the audit trap}\label{ch15:sec:trap}

\texttt{\#print axioms foo} makes Lean walk the proof term of \lean{foo}, and of every declaration
that term mentions, and list the axioms reached. It is the only reliable way to know whether a
theorem depends on \lean{sorryAx}, the axiom that a \lean{sorry} silently
introduces (or, while it was an axiom, on \lean{VOA.pbw}).

There is a trap, recorded in the project's working rules and in \leanpath{README.md}. The
package's \leanpath{lakefile.toml} sets an elaboration option,
\begin{plaincode}
[leanOptions]
pp.unicode.fun = true # pretty-prints `fun a ↦ b`
relaxedAutoImplicit = false
weak.linter.mathlibStandardSet = true
maxSynthPendingDepth = 3
\end{plaincode}
which raises the depth to which type-class synthesis may postpone sub-problems. \lean{lake build}
passes this option to every file it compiles. A bare \lean{lake env lean file.lean} does
\emph{not}: it runs Lean in the package's environment but with Lean's default options. With the
default depth, some instance problems in the library fail; Lean then reports an elaboration error
and, to continue, replaces the failed term by \lean{sorryAx}. A \texttt{\#print axioms} run in that
mode can therefore report \lean{sorryAx} for a theorem that is in fact complete, or, worse, a file
of probes can appear to succeed while auditing something other than what the build produced. The
rule is: audit only through \lean{lake build}, or through
\begin{plaincode}
lake env lean -DmaxSynthPendingDepth=3 probe.lean
\end{plaincode}
which supplies the option explicitly. The maintained audit script does the latter for its newer
probes.

\subsection{Unconditional versus conditional}\label{ch15:sec:uncond}

\texttt{\#print axioms} does not list hypotheses. The theorem
\begin{leancode}
theorem moonshine_of_frameLocality :
    ∃ x : ZMod 2 → TP, (∀ s, x s ∈ TPframe s) ∧
      ((∀ i j, TPBlockLocal leechWbar (x i) (x j)) →
        ∃ V : VertexOperatorAlgebra Vnat,
          V.Y = Ynat wbarL ∧ V.vacuum = vacNat ∧ V.conformalVector = omegaNat ∧
          V.centralCharge = 24 ∧
          V.grading 1 = ⊥ ∧
          finrank ℂ (V.grading 2) = 196884 ∧
          (∀ n : ℕ, (finrank ℂ (V.grading (n : ℤ)) : ℂ) = PowerSeries.coeff n JShift) ∧
          V.IsSimple) := by
\end{leancode}
(\leanpath{VOA/Moonshine/Main.lean}) has no binders at all, and its axiom list is clean; yet it is
\emph{conditional}: inside the statement, after the existential, sits an implication whose premise
\lean{∀ i j, TPBlockLocal leechWbar (x i) (x j)} (four instances of the twisted-sector locality
relation (4.23) of Dolan--Goddard--Montague, Chapter~\ref{ch:12}) was, at the time, not proved.
The project's rule of language, stated in \leanpath{README.md} and in the project's working rules, is:

\begin{quote}
Say ``\texttt{\#print axioms}-clean, conditional on \ldots'' for a declaration whose proof term uses
only the standard axioms but whose statement carries a hypothesis. Reserve \emph{unconditional} for a
theorem whose full hypothesis closure contains no unestablished datum.
\end{quote}

The \emph{hypothesis closure} is found by tracing every hypothesis, including those hidden in a
statement as premises of implications and those carried as fields of structures, down to its
leaves. The \leanpath{README.md} adds two further audit items for every headline result:
\emph{conclusion strength} (is the conclusion the intended one, or a weaker variant such as
``there exists some datum with property $P$'' where a pinned datum was intended?) and
\emph{satisfiability of carried data} (can the hypotheses be met at all?). The Moonshine
construction illustrates all three. Its final theorem is unconditional in the strict sense: no
binder, no premise inside the statement. Its conclusion had to be strengthened once (simplicity was
added to exclude the square-zero extension, Section~\ref{ch15:sec:trap-sqzero}). And several
intermediate reductions had to be guarded against vacuous satisfaction by trivial data (the
\lean{hits_cyclic} clause, Section~\ref{ch15:sec:phi0}).

\subsection{Independent re-verification of agent-produced work}\label{ch15:sec:agents}

Most of the roughly 155 files of \leanpath{VOA/Moonshine/} were written by background agents (Lean
sessions running in parallel, each assigned one milestone and allowed to write only new files),
coordinated by a supervising session that followed the operating rules recorded in
\leanpath{Notes/MonsterRoadmap.md}. Those rules make agent output count only after an independent
re-check by the supervisor:
\begin{enumerate}
\item read the roadmap and the working agreements at every iteration;
\item for every finished agent, build its \emph{new} files in isolation
  (\lean{lake build VOA.Moonshine.<File>}); grep them for \lean{sorry}, \lean{admit},
  \lean{native_decide} and \lean{axiom}; run \texttt{\#print axioms} on the headline declarations
  via \lean{lake env lean -DmaxSynthPendingDepth=3};
\item only then wire the files into \leanpath{VOA/Moonshine.lean};
\item update the status table with exact declaration names, and flip a status marker only after a
  green aggregate build;
\item never weaken a statement, never introduce \lean{sorry} or an axiom, never commit or push (the
  maintainers' decision), never run two builds concurrently in one tree, never touch unrelated work
  (the parallel $\mathfrak{sl}_3$ thread);
\item periodically run \lean{scripts/check_repository.py} and an aggregate build, and at
  checkpoints write a dated hand-off record (commit, build state, open residuals, dead ends) for the next session.
\end{enumerate}
The agents' own reports were treated as claims, not facts. Concretely, this caught: unwired drafts
(listed by the reachability check); a draft with \lean{sorry} (Section~\ref{ch15:sec:checkrepo});
agent notes that described their files as ``not wired'' long after wiring, or a residual as open
after it had been closed (the route notes carry ``Update'' sections and the gluing plan carries
dated ``Status and corrections'' sections for exactly this reason); and, most importantly, several
mathematical errors in the plan itself, recorded as numbered ``Corrections'' in
\leanpath{Notes/M8GluingPlan.md} and catalogued below. A general lesson of the library's history,
learned twice, applies with full force: \emph{prose drifts behind theorems}. A
docstring saying ``open'' or ``not possible'' was false when read at least seven times in earlier
sessions; the remedy is to search for the declaration the prose claims is missing before believing
it.

\subsection{The final verification record}\label{ch15:sec:record}

\textbf{The release commit.} This edition documents the release of the library named in the front matter, at commit
\texttt{\bookcommit}. The protocol of this section --- the aggregate \lean{lake build}, \lean{scripts/check_repository.py},
\lean{scripts/audit/run_axiom_audits.py}, and \lean{sync.py check} for the book's quotations --- is run at that commit by the
release procedure. The records below were made at earlier states of the library, are dated, and say what was run.

\begin{status}[title={The release build (3 November 2026)}]
Run on the development tree at the feature freeze, the source of this release (the released tree omits the held
material listed in Chapter~\ref{ch:01}):
\begin{itemize}
\item aggregate \lean{lake build} of \lean{VOA}: \lean{Build completed successfully (20331 jobs)}, exit code 0;
\item \lean{scripts/check_repository.py}: 11{,}859 library files, 0 axiom declarations, no \lean{sorry}, \lean{admit},
  \lean{native_decide}, \lean{stop} or \lean{implemented_by}; all files reachable from the umbrella except four held
  files that the release does not contain;
\item \lean{scripts/audit/run_axiom_audits.py}: passed, 5{,}958 reports across 24 probes;
\item \texttt{\#print axioms} on the 28 headline declarations of the release (including
  \lean{VOA.Moonshine.moonshine}, \lean{A2_classicallyFree_final}, \lean{D4_classicallyFree_of_inputT} and
  \lean{twoSeries_gradedUnitary_of_affine_final}): every one reports exactly
  \lean{[propext, Classical.choice, Quot.sound]}.
\end{itemize}
\end{status}


\begin{status}
Recorded at the end of the Moonshine construction (the project's checkpoint record of 7 October 2026),
on the tree that became the Moonshine commit (``Construct the Moonshine module $V^\natural$ (FLM)
unconditionally''), then the head of the main branch:
\begin{itemize}
\item aggregate \lean{lake build}: \lean{Build completed successfully (10343 jobs)}, exit code 0;
\item \lean{scripts/check_repository.py}: passed, 1867 library files, all reachable, one sanctioned
  axiom declaration (\lean{VOA.pbw}, then an axiom; 0 since it became a theorem);
\item \lean{scripts/audit/run_axiom_audits.py}: passed, 4876 reports across 20 probes (which do not
  cover \lean{VOA.Moonshine});
\item \texttt{\#print axioms VOA.Moonshine.moonshine}, via \lean{import VOA} with
  \lean{-DmaxSynthPendingDepth=3}: \lean{[propext, Classical.choice, Quot.sound]}; no \lean{VOA.pbw};
\item the core definitions (\leanpath{VOA/Core}, \leanpath{VOA/Foundations},
  \leanpath{VOA/Modularity/JFunction.lean}, \leanpath{VOA/Modules/Fusion.lean}) are unchanged since
  the parent of the Moonshine commit; the only change to a pre-existing definition is the
  nine-line repair of \lean{TwistedModule.jacobi} in \leanpath{VOA/Modularity/Twisted.lean}
  (Section~\ref{ch15:sec:twisted}).
\end{itemize}
Re-run for this chapter on the same tree: \lean{scripts/check_repository.py} passed with
``1867 library files scanned; 1867 reachable through 4676 local import edges; 1 sanctioned axiom
declaration(s)''; \texttt{\#print axioms} gave \lean{[propext, Classical.choice, Quot.sound]} for
\lean{VOA.Moonshine.moonshine}, \lean{VOA.Moonshine.topPairs_one}, \lean{VOA.jacobiAllIndices_eq_zero}
and \lean{VOA.Moonshine.TrialityTw.golayE_not_determined}, plus the library-wide list of
Section~\ref{ch15:sec:trust}. \lean{scripts/audit/run_axiom_audits.py} was re-run and printed twenty
\lean{PASS} lines (from ``\lean{PASS RigidityProgress: 3932 declarations match the trust boundary}''
to ``\lean{PASS PairClosure: 54 declarations match the trust boundary}'') and
``Axiom audit enforcement passed: 4876 reports across 20 probes.'' (1 minute 34 seconds). The
aggregate \lean{lake build} was \emph{not} re-run for this chapter; its figure of 10343 jobs is the
one recorded at the time.
\end{status}

For the reader's convenience we reprint the headline statement once more
(\leanpath{VOA/Moonshine/TopPairs.lean}):
\begin{leancode}
theorem moonshine :
    ∃ V : VertexOperatorAlgebra Vnat,
      V.Y = Ynat VnatAxioms.wbarL ∧ V.vacuum = vacNat ∧ V.conformalVector = omegaNat ∧
      V.centralCharge = 24 ∧
      V.grading 1 = ⊥ ∧
      Module.finrank ℂ (V.grading 2) = 196884 ∧
      (∀ n : ℕ, (Module.finrank ℂ (V.grading (n : ℤ)) : ℂ) = PowerSeries.coeff n JShift) ∧
      V.IsSimple :=
  moonshine_of_topPairs (by norm_num) topPairs_one
\end{leancode}

\begin{dictionary}
\textbf{Fidelity check of the headline, from the verification point of view.} The full fidelity
check, with every term traced to its definition, is in Chapter~\ref{ch:14}; here we record only the
points that the protocol of this section is designed to protect.
\begin{itemize}
\item \emph{No hypotheses.} The statement has no binders and no premise. Its proof is the
  composition of the conditional theorem \lean{moonshine_of_topPairs} (whose premises are
  $\varepsilon^2=1$, discharged by \lean{norm_num} at $\varepsilon=1$, and the finite identity
  \lean{TopPairs}) with \lean{topPairs_one}, which proves that identity. So the hypothesis closure is
  empty.
\item \emph{The existential is pinned.} ``$\exists V$'' alone would be weak: some VOA structure on
  the vector space \lean{Vnat}. The four equations \lean{V.Y = Ynat VnatAxioms.wbarL},
  \lean{V.vacuum = vacNat}, \lean{V.conformalVector = omegaNat} pin the vertex operator map to the
  explicit block map of Frenkel--Lepowsky--Meurman, built from the untwisted fields, the skew
  intertwiner $W$ and its adjoint $\overline W$ (Chapters~\ref{ch:09}--\ref{ch:10}); the vacuum and
  conformal vector are pinned likewise.
\item \emph{The non-triviality clause.} \lean{V.IsSimple} excludes the square-zero extension, which
  satisfies every other clause (Section~\ref{ch15:sec:trap-sqzero}). In addition \lean{wbarL} is the
  genuine adjoint, proved nonzero by \lean{leechWbar_ne_zero}.
\item \emph{Graded dimension.} \lean{JShift} is the pre-existing series $E_4^3/(\Delta/q)-744q=qJ(q)$
  of \leanpath{VOA/Modularity/JFunction.lean}, untouched by the construction, so the identity
  $\dim V_n=[q^n]\,qJ(q)$ is tested against an independent definition of $J$.
\item \emph{Axioms.} Standard only; no \lean{VOA.pbw}.
\end{itemize}
What the statement does \emph{not} say: nothing about $\Aut(V)$, the Monster, or uniqueness;
see Section~\ref{ch15:sec:open}.
\end{dictionary}

\begin{physics}
The verified statement underwrites the physicist's description of the Moonshine conformal field
theory: a chiral $c=24$ CFT (holomorphy, i.e.\ a single sector, is the physical expectation and is
not proved in Lean) obtained by orbifolding 24 chiral bosons
compactified on the Leech torus by the reflection $X\mapsto -X$, keeping the invariant untwisted
states and the invariant twisted states (twist fields of weight $3/2$ dressed by oscillators). The
clause $V_1=0$ says the theory has no conserved spin-one currents, hence no continuous symmetry;
$\dim V_2=196884$ counts the stress tensor plus the $196883$ primaries of the Griess algebra; the
graded dimension is the torus partition function $Z(\tau)=J(\tau)$ up to the factor $q$; uniqueness of
the vacuum (the only state of weight $0$ up to scale) is CFT type, built into the library's definition
of a vertex operator algebra; simplicity says that no consistent subsector decouples, i.e.\ the operator product expansion genuinely couples the twisted states
back into the untwisted ones ($\overline W\neq 0$).
\end{physics}

\subsection{The book's own verification record (October 2026)}\label{ch15:sec:bookrecord}

The protocol above certifies the \emph{library}. The book needs a second, separate record: that what it prints is what
the library says. By the book's rules (its constitution, Article~9), the book documents the library as of the commit
stored in \leanpath{Notes/Textbook/manifest.json}, which records, for every \lean{.lean} file, the chapter that
documents it and the hash of its contents at that commit. The script \leanpath{scripts/textbook/sync.py} then checks the
chapters against the sources. Its \lean{check} command reports, among other things, four kinds of failure that are
treated as blocking:
\begin{itemize}
\item \textbf{STALE QUOTE}: a \lean{leancode} block that no longer occurs verbatim in the source. This breaks the book's
  central promise, that the Lean text the referee reads is the Lean text the kernel checked.
\item \textbf{MISSING PATH}: a quoted declaration with no nearby \texttt{\textbackslash leanpath} citation of the file that defines it, so
  the reader cannot find the source.
\item \textbf{NO PHYSICS}: a definition or theorem not followed by a physical-meaning box, which the constitution requires
  for every statement.
\item \textbf{UNASSIGNED} files (reported by \lean{sync.py coverage}): a library file that no chapter is responsible for.
\end{itemize}
These checks are textual. They confirm that the quotations are verbatim and attributed, not that the quoted files
compile together; that is the job of the aggregate build of Section~\ref{ch15:sec:build}, and the two must not be
confused.

\begin{status}
The book's record at the end of the October 2026 update, as recorded in
\leanpath{Notes/Textbook/BOOK\_CAMPAIGN.md} (``Status'') and in the commit messages of the book:
\begin{itemize}
\item \textbf{Manifest snapshot} at a pre-release commit, with \lean{sync.py coverage} reporting $0$ unassigned
  files. Every \lean{.lean} file of the library is therefore assigned to exactly one chapter.
\item \textbf{\lean{sync.py check}} at that snapshot: $0$ STALE QUOTE, $0$ MISSING PATH, $0$ NO PHYSICS.
\item \textbf{Full book} at that snapshot: \lean{book.pdf}, $2118$ pages, $0$ undefined references. (The release
  edition, \lean{main_release.pdf}, then had $2201$ pages.)
\item \textbf{Brief}: \lean{brief.pdf}, $199$ pages; all $1199$ of its $1199$ Lean chunks verbatim against the sources at
  that snapshot.
\item \textbf{Chapters in preparation.} The chapter on affine $\mathfrak{sl}_3$ is not part of this edition. The two
  chapters on graded unitarity through the BRST coset, then held, are part of the v1.0 release edition
  (Chapters~\ref{ch:L28} and~\ref{ch:L29}); their verification records are in those chapters. Chapter~\ref{ch:L27},
  held at this snapshot, was rewritten around the $A_2$ theorem without hypotheses and is included.
\end{itemize}
\textbf{Not done at that snapshot.} The book was \emph{not} checked against an aggregate \lean{lake build} at that snapshot. The
last aggregate build recorded green is that of 26~October~2026, at the $A_2$ milestone of the Deligne campaign:
$13868$ jobs, with the axiom audit passing (the project's build record; not re-run for the book). By the project's policy the aggregate build is run
only before a merge to the main branch, and it is memory-heavy, so it was not re-run for the book. Since then, the aggregate build \texttt{bash scripts/campaign\_build.sh VOA} (the full \texttt{lake build VOA} under the project's memory-capped wrapper) was green in two pre-release builds on 2026-10-28 (18{,}901 jobs, and 18{,}907 jobs after merging the then-current \texttt{origin/main}); \texttt{scripts/check\_repository.py} passed (10{,}431 library files, all reachable, 0 project axioms) and \texttt{scripts/audit/run\_axiom\_audits.py} passed (5{,}958 reports across 24 probes). For the release commit see Section~\ref{ch15:sec:record}. The quotations are
verbatim against the sources at the manifest commit; their compilation at that commit is inferred from the build of
26~October and from the targeted builds recorded in the project's internal build records, not re-confirmed.
\end{status}

\paragraph{The trust base at the manifest commit.} The library has \emph{no} project axiom. \lean{VOA.pbw}, which until
19 October 2026 was the single sanctioned axiom (Section~\ref{ch15:sec:trust}), is now the theorem
\lean{theorem pbw : PBWPrinciple} in \leanpath{VOA/Foundations/AffineLie.lean}, proved through the files
\leanpath{VOA/Foundations/PBWTheorem/} (Chapter~\ref{ch:L02}). \texttt{\#print axioms VOA.pbw} lists only Lean's three
standard axioms; this was re-run for Chapter~\ref{ch:L02} and again, independently, by the release audit of
Chapter~\ref{ch:L22} (Section~\ref{ch15:sec:releaseaudits}). Two of the release audits also searched \leanpath{VOA/} for
top-level \lean{axiom} declarations and found none: every occurrence of the word is in a comment. Consequently the
register ``proved using the project axiom'' of the introduction is empty, and ``\texttt{\#print axioms}-clean'' now means,
without exception, that the axiom list is a subset of \lean{[propext, Classical.choice, Quot.sound]}. What it still does
\emph{not} mean is ``unconditional'' (Section~\ref{ch15:sec:uncond}): the new chapters contain many
\texttt{\#print axioms}-clean theorems that carry open named hypotheses, and they are listed in
Section~\ref{ch15:sec:open-later}.

\paragraph{Two further facts of the record.}
\begin{itemize}
\item \textbf{The repository check passes.} \lean{scripts/check_repository.py} (Section~\ref{ch15:sec:checkrepo})
  failed for a time on 27~October~2026 because of unwired files; after they were wired in or moved, the Foundations
  release audit recorded it passing ($7{,}839$ files, all reachable). Re-run on the working tree of the branch
  \lean{deligne-lean} on 28~October~2026, it reported: ``Repository check passed: 10353 library files scanned; 10353
  reachable through 24571 local import edges; 0 sanctioned axiom declaration(s)''. That working tree had uncommitted
  changes; for the release commit see Section~\ref{ch15:sec:record}.
\item \textbf{Docstrings were edited after the last aggregate build.} A first batch of stale docstrings (45 files,
  comments only) was corrected on 27~October; a targeted build of those 45 modules was green ($13{,}276$ jobs). The
  release audits below found further edits of the same kind in the import closures they examined and inspected them as
  comment-only. A comment cannot change a theorem, but an edited file must still be rebuilt, which is one more reason
  why every audit asks for an aggregate build at the release commit.
\end{itemize}

\begin{physics}
The physicist's analogue is the difference between a calculation and its write-up. The library is the calculation; the
book is the paper that reports it. A paper can misquote its own calculation (a sign dropped in transcription, a
hypothesis forgotten in the abstract), and the sync check is the systematic comparison of the two. It does not re-run the
calculation; the aggregate build does that.
\end{physics}

\subsection{The independent release audits of 27--28 October 2026}\label{ch15:sec:releaseaudits}

Before a public release the project decided (\leanpath{Notes/RELEASE\_PLAN.md}, ``Release tiers'') that every headline
theorem to be named in an announcement must pass one independent \emph{release audit}: a fresh referee, who had not
written the chapter, compares the statement with the source paper, traces every definition it depends on down to its
leaves, lists every explicit hypothesis, and re-runs \texttt{\#print axioms}. The chapter is used only as a map; every
conclusion is checked against Lean. The project's record of over-claims caught late (three in the week before the
audits alone) is the reason this step exists. Twelve such audits were run on 27 and 28~October~2026; their verdicts
are recorded in \leanpath{Notes/RELEASE\_PLAN.md} (``Release
audit results''). One line each:
\begin{itemize}
\item \lean{RELEASE_AUDIT_L16.md} (Chapter~\ref{ch:L16}, Arakawa--Kawasetsu), 27~October: the proposed announcement is not
  announceable as worded; 1~BLOCKING (the wording), 4~MAJOR, 6~MINOR. Discussed below.
\item \lean{RELEASE_AUDIT_L20.md} (Chapter~\ref{ch:L20}, Arakawa--Dai), 27~October: no blocker; 0~BLOCKING, 1~MAJOR (the
  release build), 7~MINOR. Discussed below.
\item \lean{RELEASE_AUDIT_L21.md} (Chapter~\ref{ch:L21}, Huang and related), 27~October: no blocker; 0~BLOCKING, 3~MAJOR,
  8~MINOR. Discussed below.
\item \lean{RELEASE_AUDIT_L22.md} (Chapter~\ref{ch:L22}, Gannon--Rayhaun), 27~October: announceable with caveats;
  0~BLOCKING, 2~MAJOR, 5~MINOR. Discussed below.
\item \lean{RELEASE_AUDIT_MOONSHINE.md} (Chapters~\ref{ch:02}--\ref{ch:14}, $\Vnat$), 27~October: announceable with caveats,
  wording only; 0~BLOCKING, 0~MAJOR, 6~MINOR. \lean{moonshine} takes no hypotheses, has standard axioms and a unique
  witness; $J$ is a formal $q$-series; the isomorphism with FLM's explicit $\Vnat$ is on paper only; no Monster,
  holomorphy or uniqueness claim.
\item \lean{RELEASE_AUDIT_RIGIDITY.md} (Chapters~\ref{ch:L07} and~\ref{ch:L09}), 27~October: announceable with caveats;
  0~BLOCKING, 0~MAJOR, 3~MINOR (say ``graded first-order''; non-vacuity argued on paper only).
\item \lean{RELEASE_AUDIT_ZHU.md} (Chapters~\ref{ch:L12}--\ref{ch:L15}), 27~October: no headline blocked, but ``Zhu's
  modularity theorem is formalized'' is misleading without qualification; 3~HIGH, 4~MEDIUM, 6~lower. Its finding that
  ``rational $\Rightarrow$ $A(V)$ semisimple'' is unproved was superseded the next day (see the new-declarations audit).
\item \lean{RELEASE_AUDIT_FOUNDATIONS.md} (core definitions and constructions), 27~October: faithful with caveats;
  0~BLOCKING, 2~MAJOR, 7~MINOR (every headline must say ``of CFT type''; the structure now called \lean{PreVertexAlgebra} lacks translation
  covariance, Section~\ref{ch15:sec:debt}).
\item \lean{RELEASE_AUDIT_L17_L18_L23.md} (Chapters~\ref{ch:L17}, \ref{ch:L18}, \ref{ch:L23}), 27~October: 0~BLOCKING,
  7~MAJOR, 14~MINOR; includes the \lean{SuperTwistedModule} defect of Section~\ref{ch15:sec:super}.
\item \lean{RELEASE_AUDIT_L24_L25_A2.md} (Chapters~\ref{ch:L24}, \ref{ch:L25} and the $A_2$ final theorem), 28~October:
  1~HIGH, 5~MEDIUM, 8~LOW; the $A_2$ statement is faithful, its announcement blocked until green umbrella and aggregate
  builds exist at the release commit (condition met on 28~October: Chapter~\ref{ch:L27}, Section~\ref{l27:sec:verification}).
\item \lean{RELEASE_AUDIT_NEW_DECLS.md} (the declarations added on 28~October), 28~October: all announceable with
  caveats; 0~BLOCKING, 4~MEDIUM, 7~LOW.
\item \lean{RELEASE_AUDIT_FOLLOWUPS.md} (the other account's follow-ups), 28~October: none wrong, none blocked; 2
  announceable as stated, 4 with caveats.
\end{itemize}
A separate precision audit of the Deligne literature inputs (\lean{DELIGNE_INPUTS_PRECISION.md}, 28~October) found the
total-dimension bound and the Demazure realizations published for all six types and the star bound unpublished for
$G_2$, $F_4$ and the $E$-types. The paragraphs below record the first four audits in detail, as reported on 27~October;
where Lean results of 28~October have since removed a caveat, this is said in the paragraph. The corrections the
audits asked of their chapters have been applied.

\paragraph{What the auditors of 27~October ran, and what they did not.} No auditor ran \lean{lake build}. Each wrote scratch files
importing the headline modules and ran them with \lean{lake env lean -DmaxSynthPendingDepth=3} (the safe form of
Section~\ref{ch15:sec:trap}) against the compiled \lean{.olean} files already present, dated 25--26~October; the cost was
2.6--9.5 seconds and 3.4--8.1\,GB per file. Each audit was made at a different pre-release commit of a working tree with uncommitted
changes. Each
auditor checked that the uncommitted changes in the relevant import closure were comment-only, or that the closure was
unchanged. \emph{Every one of the four reports therefore carries the same process caveat}: the aggregate build, the
repository check and the axiom audit must be re-run, green, on the committed release commit before anything is
announced. All three were green on 28~October~2026 (paragraph \emph{Not done at that snapshot} above); for the release
commit see Section~\ref{ch15:sec:record}.

\paragraph{Chapter~\ref{ch:L16} (Arakawa--Kawasetsu): the proposed announcement is not announceable as worded.} Counts:
1~BLOCKING, 4~MAJOR, 6~MINOR. The BLOCKING finding concerns the wording of the proposed announcement, not the Lean code.
That wording (``the Arakawa--Kawasetsu theorem, characters of quasi-lisse VOAs satisfy an MLDE, and an explicit
non-lisse quasi-lisse example are formalized'') is false on three counts. (1)~\emph{No Lean theorem proves the analytic Arakawa--Kawasetsu
MLDE for any vertex algebra.} (The audit's wording, ``no MLDE for any VOA'', was slightly too broad: Chapter~\ref{ch:L25}
proves formal, coefficientwise second-order MLDEs for the vacuum characters of the seven Deligne-series VOAs.) Within
Chapter~\ref{ch:L16}, the only theorem with a genuine MLDE conclusion,
\lean{VOA.IsSqPrimary.exists_monicMLDEOn_all_analyticBlocks_of_etingofFinite}, is \texttt{\#print axioms}-clean but
conditional on two hypotheses that were inhabited nowhere in the library, \lean{AKWeightGrading W} and
\lean{E46SerreClosed} (Section~\ref{ch15:sec:open-later}). Since 28~October the second is a theorem,
\lean{VOA.e46SerreClosed} (\leanpath{VOA/Modularity/EisensteinSerreClosure.lean}), and since 29~October the first is a theorem for every
VOA of the library (CFT type), \lean{VOA.akWeightGrading} (\leanpath{VOA/Modularity/AKGradingTotal.lean}). The MLDE is
therefore conditional only on Etingof-finiteness:
\lean{VOA.IsSqPrimary.exists_monicMLDEOn_all_analyticBlocks_of_etingofFinite_uncond} (despite the suffix, it keeps
\lean{EtingofFinite W}) gives one monic MLDE of some unbounded order $s\ge1$, coefficients in $\C[E_4,E_6]$ of weight
$2(s-i)$, satisfied by $F(a)$ for every abstract analytic genus-one block $F$; and
\lean{VOA.exists_monicMLDEOn_all_analyticBlocks_simpleSl2QL} is its unconditional instance at the vacuum of
$L_{-4/3}(\mathfrak{sl}_2)$, of unspecified order (Arakawa--Kawasetsu's order~2 is not proved; that the vacuum character
is such a block follows from library lemmas but is not packaged as one theorem). Both are \texttt{\#print axioms}-clean
(recorded in \texttt{Notes/release/WS-RELEASE2.md}); the aggregate build at the release commit is pending. (2)~\emph{Quasi-lisseness is not formalized}: there is no
associated variety and no symplectic leaf; the theorems assume \lean{EtingofFinite}, finite-dimensionality of the zeroth
Poisson homology $R_V/\{R_V,R_V\}$, which is the \emph{conclusion} of the Etingof--Schedler theorem (logically more
general, but not ``quasi-lisse''). (3)~\emph{The unconditional ``character ODE'' theorems} at $L_{-4/3}(\mathfrak{sl}_2)$
(\lean{character_formalODE_simpleSl2QL_uncond} and siblings) have arbitrary subexponential power series as
coefficients, chosen after the module; they are near-vacuous growth statements, not MLDEs, although a docstring still
calls one of them ``the Arakawa--Kawasetsu character equation'' (Chapter~\ref{ch:L16}, Section~\ref{l16:sec:weak}).
A further MAJOR: the Zhu-modularity chain under Etingof-finiteness has its rank hypothesis discharged only for the trivial
VOA, and it is expected to be false at $L_{-4/3}(\mathfrak{sl}_2)$ (a paper argument). \emph{What is announceable}, with
standard axioms only: $L_{-4/3}(\mathfrak{sl}_2)$ is not $C_2$-cofinite yet has finite-dimensional zeroth Poisson
homology (\lean{not_C2Cofinite_simpleSl2QL}, \lean{etingofFinite_simpleSl2QL'},
\lean{strictEtingofFinite_simpleSl2QL}); for such VOAs characters of ordinary modules converge on the upper half-plane
and the genus-one conformal blocks form a finite-dimensional space; and $C_2$-cofinite implies Etingof-finite. The
MLDE may be mentioned only as formalized conditionally on \lean{EtingofFinite} (the stand-in for quasi-lisseness), with an
unspecified order, and unconditionally at $L_{-4/3}(\mathfrak{sl}_2)$ in that form. Consequence recorded in the release
plan: Chapter~\ref{ch:L16} is released, but it leaves the announcement's list of headlines.

\paragraph{Chapter~\ref{ch:L20} (Arakawa--Dai): no blocker.} Counts: 0~BLOCKING, 1~MAJOR, 7~MINOR. All ten headline
declarations of \leanpath{VOA/Examples/AffineSl2BoundaryMainTheorems.lean} are \texttt{\#print axioms}-clean, with the single
hypothesis $1\le m$; the definitions trace to the genuine objects (the actual simple quotient, the real Li filtration, a
genuine Poisson vertex algebra, the character computed from actual dimensions). The four character theorems
(\lean{coeff_product_le_character}, \lean{character_eq_product}, \lean{character_eq_chain},
\lean{character_eq_fermionic_chain}, the last being the paper's Theorem~1.2 display) and
\lean{arcPoissonEquiv_conformal_iff} are announceable as stated. The caveats: \lean{canonicalMap_bijective},
\lean{boundary_classically_free} and \lean{arcPoissonEquiv} are stated on the arc algebra of the \emph{presented} algebra
$\C[e,h,f]/(eQ^m,hQ^m,fQ^m)$, and no single Lean theorem has the paper's shape ``the canonical map
$J_\infty R_V\to\operatorname{gr}^FV$ is bijective''; the paper's Theorem~1.1 follows from two Lean theorems together
(the presented isomorphism and \lean{baseC2AlgEquiv}) and should be announced in that two-part form (audit verdict of
27~October; superseded on 28~October, see below and Chapter~\ref{ch:L20}, Theorem~\ref{l20:thm:ad11}; naturality proved).
\lean{actualC2ArcEquiv} is a bare bijection of hom-sets for each coefficient algebra, with no naturality stated, so it
must not be announced as a representability statement; and \lean{baseC2AlgEquiv} is canonical by its definition, not by
its type. The single MAJOR is the release build (above). Chapter~\ref{ch:L20}, Section~\ref{l20:sec:defects}, records the
same points. \emph{Since 28~October} two of these caveats are removed in Lean, in the namespace
\lean{VOA.AffineSl2BoundaryClassicalFreeness}: \lean{arakawaDai_theorem_1_1}
(\leanpath{VOA/Examples/AffineSl2BoundaryTheorem11.lean}) states the paper's Theorem~1.1 as one theorem, and
\lean{actualC2ArcEquiv_natural} (\leanpath{VOA/Examples/AffineSl2BoundaryMainTheorems.lean}) states the naturality, so
representability holds; the follow-ups audit rated both ``as stated''. What remains is that $J_\infty R_V$ is realized
as \lean{PresentedJet} with a universal arc, since the library has no $J_\infty$ functor on abstract algebras.

\paragraph{Chapter~\ref{ch:L21} (Huang and related): no blocker.} Counts: 0~BLOCKING, 3~MAJOR, 8~MINOR. Every headline is
\texttt{\#print axioms}-clean and instantiated at the lattice VOA $V_2$, so none is vacuous. Verdicts:
\begin{itemize}
\item the braided monoidal category of ordinary modules of a strongly rational VOA (a genuine Mathlib
  \lean{BraidedCategory}, on the Huang--Lepowsky tensor product, with pentagon and hexagons proved), and the logarithmic
  braided tensor category of a $C_2$-cofinite VOA: \emph{announceable as stated};
\item rigidity, Verlinde and the modular category: \emph{with caveats}. (F1, MAJOR) They are proved for \lean{Nice} VOAs,
  i.e.\ strongly rational \emph{and pseudo-unitary}, although Huang's theorems do not need pseudo-unitarity; the
  hypothesis is genuinely used in the proofs. (Since 29~October its use for $V'\cong V$ is removed:
  \lean{VOA.contra_vacuum_equiv_of_stronglyRational} and the \lean{_of_stronglyRational} copies in
  \leanpath{VOA/Huang/VerlindeVacuumFusionSR.lean} give the vacuum-fusion chain for strongly rational VOAs; the
  Cardy-positivity step \lean{vacColumn_mem_K} keeps it, essentially, and the headline theorems still take \lean{Nice}.) Rigidity was right rigidity only (\lean{rightRigid_ordinary}); since
  28~October the two-sided \lean{RigidCategory} is a theorem, \lean{VOA.Huang.HuangRigid.rigid_ordinary}
  (\leanpath{VOA/Huang/RigidOrdinary.lean}), on the same monoidal instance as the braiding (checked by the
  new-declarations audit), still for \lean{Nice} VOAs. Verlinde is proved in the form
  $S_{ax}S_{bx}=S_{0x}\sum_cN_{ab}^cS_{cx}$ with $S_{0x}\neq0$;
\item the congruence property: \emph{not announceable until it has a named theorem} (F2, MAJOR). It is true in the
  library (it is discharged inside the proof of the Gannon--Rayhaun theorem, and the auditor's scratch declarations
  \lean{audit_dln_congruence} and \lean{audit_zhuCongruence} compile with standard axioms), but every library declaration
  that concluded it carried hypotheses. Since 28~October the named theorems exist, for \lean{Nice} VOAs:
  \lean{dlnCongruenceInput_holds}, \lean{zhuCongruenceInput_holds} and \lean{zhu_kernelContainsGamma_at_level}, in
  the namespace \lean{VOA.Huang.CongruenceNamed} (\leanpath{VOA/Huang/CongruenceNamed.lean}). (F3, MAJOR, wording) Only the kernel half of Ng--Schauenburg's Theorem~6.8 is
  proved (\lean{nsProjectiveKernelAll_holds}); the half saying that the level equals the Frobenius--Schur exponent is
  not. (The new-declarations audit adds a precision: at the minimal level the kernel statement is the full kernel
  claim of Dong--Lin--Ng; ``kernel half'' refers only to the Ng--Schauenburg input.)
\item Deligne's theorem (\lean{deligneTannakian_holds}): \emph{with caveats}; it is the positive-dimension (Tannakian,
  bosonic) case for symmetric fusion categories, not the super version.
\end{itemize}

\paragraph{Chapter~\ref{ch:L22} (Gannon--Rayhaun): announceable with caveats.} Counts: 0~BLOCKING, 2~MAJOR, 5~MINOR. The
three headlines (\lean{gannonRayhaun_forward}, \lean{gannonRayhaunConverse_holds}, \lean{gannonRayhaun_classification})
have no hypotheses, their axioms are standard, and the statement matches Gannon--Rayhaun's Theorem~3.1 and
Proposition~3.2 (the same list, both directions, pairwise non-isomorphism). Non-vacuity is proved (\lean{V2m_nice},
\lean{exceptionalV2Group_exists}). The caveats: (M1, MAJOR, wording) ``rational'' is formalized as
\lean{AdmissiblyRational}: every weak (mode-stable, not necessarily graded) submodule of an admissible module has a
weak complement. The library's own note (\leanpath{VOA/Modules/AdmissibleRationality.lean}) says that its relation to
Dong--Li--Mason rationality is not settled and that neither implication is proved in Lean; the chapter overstated the
match and has been corrected (Chapter~\ref{ch:L22}, Remark~\ref{l22:rem:rational}). Since 29~October one direction is a
theorem: for $C_2$-cofinite $V$, Dong--Li--Mason rationality (formalized faithfully as \lean{DLMRational}) implies
\lean{AdmissiblyRational} (\lean{VOA.Classification.COne.DLMRat.admissiblyRational_of_dlmRational}); the converse is
open. (M2, MAJOR) the release build. The
minor points: the combined statement and the converse are stated for carriers in \lean{Type 0} only (disclosed in
Remark~\ref{l22:rem:universe}); rationality and pseudo-unitarity quantify over modules in the VOA's own universe; and
stale docstrings.

\begin{physics}
In physical language the audits separate three things that a CFT paper often runs together. ``The theory is
consistent'' (the braided tensor category exists) was confirmed as stated. ``The theory is unitary in the weak sense
needed for positivity arguments'' (pseudo-unitarity, quantum dimensions equal to the ratios of $S$-matrix entries) turns
out to be an extra input in the Lean proofs of rigidity and the Verlinde formula (now essentially only in the
Cardy-positivity step), while Huang's proofs do without it.
And ``the characters obey a modular differential equation'', the physicist's working definition of a theory with
finitely many characters, is proved with a specified order for no theory analytically: formal coefficientwise
second-order MLDEs are proved for the seven Deligne vacuum characters (Chapter~\ref{ch:L25}); the analytic
Arakawa--Kawasetsu MLDE, of unspecified order and for abstract genus-one blocks, is proved conditionally only on
Etingof-finiteness, and unconditionally at the vacuum of $L_{-4/3}(\mathfrak{sl}_2)$. The audits'
contribution is to say which of these sentences may be printed.
\end{physics}

% ---------------------------------------------------------------------------------------------
\section{A catalogue of dead ends, traps and corrected mistakes}\label{ch15:sec:deadends}

The working agreements of the project single out the record of dead ends as \emph{the
highest-value section} of every hand-off record between sessions: ``the project's worst time sinks have all been
re-attempts of routes already known to fail''. This section consolidates that record. Each entry
states (i) the claim, definition or route; (ii) why it fails, with the counterexample or
certificate; (iii) the lesson. Entries backed by a Lean theorem are presented with the full
fidelity protocol; entries backed by a numerical certificate or by a paper argument say so.

The entries fall into six kinds, and it helps to name them, because each kind is caught by a
different check:
\begin{enumerate}[label=(\Alph*)]
\item \emph{unsatisfiable definitions}: an axiom system no nonzero object satisfies (caught only by
  constructing an example, or by proving a certificate of emptiness);
\item \emph{conclusions satisfied by trivial objects}: the target theorem holds for a degenerate
  object as well as for the intended one (caught by asking ``what else satisfies this?'');
\item \emph{wrong formulas and conventions}: a sign, a normalisation or a modification that a
  source states differently or not at all (caught by numerical controls);
\item \emph{ambiguities}: a definition that does not determine what it was meant to determine
  (caught when a proof needs the missing information);
\item \emph{circular or insufficient reductions}: a reduction step that silently assumes its
  conclusion, or reduces to something no easier (caught by tracing what each step uses);
\item \emph{tooling and process traps}.
\end{enumerate}

\subsection{Unsatisfiable definitions}

\subsubsection{The twisted Jacobi identity over all indices}\label{ch15:sec:twisted}

Recall from Chapter~\ref{ch:06} that a $g$-twisted module, for an automorphism $g$ of order $N$,
has vertex operators $Y_M(u,z)=\sum_{r\in\frac1N\Z}u_{(r)}z^{-r-1}$, and that for $u$ in the
eigenspace $V^j=\{v: gv=e^{-2\pi i j/N}v\}$ only the modes $r\in -\frac jN+\Z$ occur. The library
encodes the mode $u_{(k/N)}$ as \lean{YM u k} with $k\in\Z$, so the support condition says
\lean{YM u k = 0} unless $k\equiv -j\pmod N$. The twisted Jacobi identity of Dong--Li--Mason is a
family of identities indexed by $\ell\in\Z$ and by the fractional indices $k,m$; in the source it is
asserted only for $k$ in the support coset of $u$. Before 6 October 2026 the field
\lean{TwistedModule.jacobi} asserted it for \emph{every} $k$. That predicate is preserved, under its
own name, as the subject of a certificate. File \leanpath{VOA/Moonshine/HeisenbergTwistedModuleDefect.lean}:
\begin{leancode}
def JacobiAllIndices (g : Aut W) (N : ℕ) {M : Type*} [AddCommGroup M] [Module ℂ M]
    (YM : 𝒱 →ₗ[ℂ] (ℤ → End ℂ M)) : Prop :=
  ∀ (j : ℕ) (u : 𝒱), u ∈ g.twistedEigenspace N j → ∀ (v : 𝒱) (x : M) (ℓ k m : ℤ),
    (∑ᶠ i : ℤ, (negOnePow i * cbz ℓ i) •
        (YM u (ℓ * N + k - i * N) (YM v (m + i * N) x)
          - negOnePow ℓ • YM v (ℓ * N + m - i * N) (YM u (k + i * N) x)))
      = ∑ᶠ i : ℕ, cbinom ((k : ℂ) / N) i • YM (W.smul u (ℓ + i) v) (k + m - i * N) x
\end{leancode}

\begin{theorem}[The all-indices twisted Jacobi identity forces $M=0$]\label{ch15:thm:jacobiAll}
Let $V$ be a vertex operator algebra, $g\in\Aut V$, $N\ge 2$, and let $Y_M$ be a linear map from $V$
to fractional-mode families of endomorphisms of a vector space $M$ such that $Y_M(\mathbf 1,z)=\mathrm{id}$
(only the mode $\mathbf 1_{(-1)}$, label $k=-N$, is nonzero, and it is the identity). If $Y_M$
satisfies the twisted Jacobi identity for all $j,u\in V^j,v,x,\ell,k,m$ without restriction on $k$,
then $M=0$.
\end{theorem}

File \leanpath{VOA/Moonshine/HeisenbergTwistedModuleDefect.lean}:
\begin{leancode}
theorem jacobiAllIndices_eq_zero {g : Aut W} {N : ℕ} (hN : 2 ≤ N) {M : Type*}
    [AddCommGroup M] [Module ℂ M] (YM : 𝒱 →ₗ[ℂ] (ℤ → End ℂ M))
    (hvac : ∀ k : ℤ, YM W.vacuum k = if k = -(N : ℤ) then 1 else 0)
    (hJ : JacobiAllIndices g N YM) (x : M) : x = 0 := by
\end{leancode}

\begin{dictionary}
\begin{itemize}
\item \lean{YM : 𝒱 →ₗ[ℂ] (ℤ → End ℂ M)}: a $\C$-linear map assigning to each state $u$ the family
  of its fractional modes, \lean{YM u k}$\,=u_{(k/N)}$. \lean{End ℂ M} is the algebra of linear
  endomorphisms of $M$.
\item \lean{hvac}: the vacuum acts by $\mathbf 1_{(-1)}=\mathrm{id}$, which in the label convention
  $k/N=-1$ is $k=-N$; all other modes vanish. This is the twisted-module vacuum axiom
  \lean{vacuum_YM}.
\item \lean{cbz ℓ i} is the integer binomial $\binom{\ell}{i}$ extended to $i\in\Z$ (zero for
  $i<0$), \lean{negOnePow i} is $(-1)^i$, and \lean{cbinom x i} is the generalised binomial
  $x(x-1)\cdots(x-i+1)/i!$ with complex upper index; \lean{W.smul u n v} is the $n$-th product
  $u_{(n)}v$ in the VOA; \lean{∑ᶠ} is a sum of finitely many nonzero terms.
\item The conclusion \lean{x = 0} for an arbitrary \lean{x : M} is $M=0$.
\end{itemize}
The theorem is \emph{stronger} than needed to condemn the old definition: it uses only the vacuum
axiom, not the other fields of a twisted module.
\end{dictionary}

\paragraph{Proof on paper.} Take $u=v=\mathbf 1$ (so $j=0$), $\ell=-1$, $k=1$, $m=-N-1$. On the left
every term contains a vacuum mode with label $-N+1-iN$ or $1+iN$; neither equals $-N$ for any
integer $i$ when $N\ge 2$ (the first would need $iN=1$, the second $(i+1)N=-1$). So the left side is
$0$. On the right the term $i$ involves $\mathbf 1_{(-1+i)}\mathbf 1$, which is $\mathbf 1$ for $i=0$
and $0$ otherwise; the $i=0$ term is $\binom{1/N}{0}\,Y_M(\mathbf 1)_{(-N)}x=x$. Hence $0=x$.
The point is that $k=1$ lies outside the support coset of $u=\mathbf 1$ (which is $k\equiv 0$), so
the identity is being asserted for modes that do not exist. The repair adds exactly the hypothesis
\texttt{(k + j) \% (N : ℤ) = 0} to the field; Chapter~\ref{ch:06} prints the repaired field and
checks that it is the Dong--Li--Mason identity.

\begin{physics}
In an orbifold CFT the twist field creates a branch cut: a field $u$ of $g$-charge $j$ picks up
the phase $e^{2\pi i j/N}$ when carried around the twisted state, so its mode expansion only
contains powers $z^{-r-1}$ with $r\in-\frac jN+\Z$. The Jacobi identity is the statement that the
operator product expansion is associative and local \emph{for the modes that exist}. Demanding it
for non-existent modes is like demanding that a fermion in the Ramond sector also have
Neveu--Schwarz modes: the only Hilbert space compatible with both is empty.
\end{physics}

\begin{status}
\lean{VOA.jacobiAllIndices_eq_zero}: unconditional (its hypotheses are exactly the data being
refuted); \texttt{\#print axioms} re-run for this chapter: \lean{[propext, Classical.choice,
Quot.sound]}. The repair is the only change to a pre-existing definition made by the Moonshine
work (diff of \leanpath{VOA/Modularity/Twisted.lean} across the Moonshine commit:
one hypothesis line, the docstring, and one binder in \lean{OrdVModule.toTwisted}, which is the
case $N=1$ and is unaffected). Before the repair, every theorem with a hypothesis
``\lean{T : TwistedModule g N M}'' with $N\ge2$ and $M\neq0$ was vacuous.
\end{status}

\begin{warning}
\textbf{Lesson.} A structure with no nontrivial instance can hide a modelling bug indefinitely
(\leanpath{README.md} calls the first instance of a definition its ``correctness oracle''). The
defect surfaced only when the first nonzero twisted module (the twisted Heisenberg module of
Chapter~\ref{ch:07}) was attempted. Every new structure should be instantiated nontrivially as
early as possible.
\end{warning}

\subsubsection{SuperTwistedModule: the same defect, not repaired}\label{ch15:sec:super}

The library's super analogue, in \leanpath{VOA/Modules/SuperTwistedModule.lean}, has the same shape.
File \leanpath{VOA/Modules/SuperTwistedModule.lean}:
\begin{leancode}
  support : ∀ (j : ℕ) (u : 𝒱), u ∈ g.twistedEigenspace N j →
    ∀ k : ℤ, (k - j) % (N : ℤ) ≠ 0 → YM u k = 0
  /-- The full Koszul twisted Jacobi identity in fractional-mode form. -/
  jacobi : ∀ {p q : ZMod 2} (j : ℕ) (u : 𝒱),
    u ∈ g.twistedEigenspace N j → u ∈ W.parity p →
    ∀ (w : 𝒱), w ∈ W.parity q → ∀ (x : M) (ℓ k m : ℤ),
    (∑ᶠ i : ℤ, (negOnePow i * cbz ℓ i) •
        (YM u (ℓ * N + k - i * N) (YM w (m + i * N) x)
          - (koszul p q * negOnePow ℓ) •
            YM w (ℓ * N + m - i * N) (YM u (k + i * N) x))) =
      ∑ᶠ i : ℕ, fractionalBinom ((k : ℂ) / N) i •
        YM (W.smul u (ℓ + i) w) (k + m - i * N) x
\end{leancode}

\begin{warning}
\textbf{Open defect.} The field \lean{jacobi} carries no restriction of $k$ to the support coset
(here written $k\equiv j\pmod N$, a different sign convention from \lean{TwistedModule}). The vacuum
is even, so the certificate of Theorem~\ref{ch15:thm:jacobiAll} applies verbatim with $p=q=0$,
\lean{koszul 0 0 = 1}: for $N\ge 2$ every \lean{SuperTwistedModule} is zero. The source file
\leanpath{VOA/Moonshine/HeisenbergTwistedModuleDefect.lean} records this in its docstring; the
super certificate itself has \emph{not} been formalised (argued on paper here). The structure is
mentioned in twelve files (\leanpath{VOA/Modules/SuperTwistedModule.lean},
\leanpath{VOA/Modules/SuperTwistedZhu.lean}, \leanpath{VOA/Modules/SuperRelativeEtingofFinite.lean},
seven files \leanpath{VOA/Modularity/Super*.lean}, \leanpath{VOA/Modularity/Twisted.lean}, and the certificate file). Any theorem among them
whose hypotheses include a \lean{SuperTwistedModule} with $N\ge2$ is vacuous until the field is
repaired. Repairing it is a decision left to the maintainers (Section~\ref{ch15:sec:open}). The
Moonshine construction does not use \lean{SuperTwistedModule}.
\end{warning}

\subsubsection{Bilinear positivity on a complex subspace}

\begin{theorem}[A $\C$-bilinear form is never positive on a nonzero $\C$-subspace]\label{ch15:thm:bilpos}
Let $\beta$ be a $\C$-bilinear form on a complex vector space $E$ and $S\le E$ a subspace. If
$\beta(x,x)$ is a positive real number for every nonzero $x\in S$, then $S=0$.
\end{theorem}

File \leanpath{VOA/Core/GradedUnitarity.lean}:
\begin{leancode}
theorem eq_bot_of_bilinear_pos (β : BilinForm ℂ E) (S : Submodule ℂ E)
    (hpos : ∀ x ∈ S, x ≠ 0 → ∃ r : ℝ, 0 < r ∧ β x x = (r : ℂ)) : S = ⊥ := by
\end{leancode}

\begin{dictionary}
\lean{BilinForm ℂ E} is a $\C$-bilinear form, \lean{Submodule ℂ E} a complex subspace,
\lean{(r : ℂ)} the real number $r$ viewed as a complex number, and \lean{S = ⊥} says $S$ is the
zero subspace. The statement is exactly Theorem~\ref{ch15:thm:bilpos}.
\end{dictionary}

\paragraph{Proof on paper.} If $x\neq0$ in $S$ has $\beta(x,x)=r>0$, then $ix\in S$ and
$\beta(ix,ix)=i^2r=-r<0$, a contradiction.

\begin{physics}
Unitarity of a CFT is positivity of a \emph{Hermitian} inner product $\langle\cdot|\cdot\rangle$
(antilinear in one slot); the Zamolodchikov metric, the BPZ form with an antilinear involution,
reflection positivity are all Hermitian. A bilinear form can only be ``positive'' on a real form of
the space. An earlier encoding of clause (vi) of the graded-unitarity conjecture of
Arabi Ardehali--Beem--Lemos--Rastelli (a positivity condition) used a bilinear form and was therefore false for every
VOA with a nonzero weight space, while its documentation called the conjecture open. Fixed by
stating clause (vi) for a sesquilinear form (\lean{SesqForm}).
\end{physics}

\begin{status}
Unconditional; standard axioms (it is part of the maintained audit). Recorded among the
library's dead ends (\leanpath{README.md}, ``Known dead ends'').
\end{status}

\subsubsection{Other unsatisfiable or refutable predicates in the library}

The library-wide register (\leanpath{README.md} ``Known dead ends'', and the project's internal session records) contains several more
entries of kind (A). We list them with their certificates; the full fidelity treatment of each is
in the chapter covering the relevant area.

\begin{itemize}
\item \textbf{\lean{GenFields} over the left ideal \lean{MRbot} is empty} (generators-and-OPE
  presentation of the Virasoro VOA). File \leanpath{VOA/Core/NonlinearVirasoro.lean}:
\begin{leancode}
theorem not_nonempty_genFields_MRbot_virasoro (c : ℂ) :
    ¬ Nonempty (GenFields (MRbot (virasoroOPE c))) := by
\end{leancode}
  The structure \lean{GenFields} asks for local fields generating a quotient of the tensor
  algebra by an ideal of relations; with the \emph{left} ideal \lean{MRbot} the creation part of
  the generating field is forced to be left multiplication, and then the mode commutator modulo a
  left ideal cannot be a sum of delta functions. The certificate is an explicit linear functional
  on words of length three that kills \lean{MRbot} but must equal $(-1)^N$ by locality of order $N$.
  A step documented as ``precisely the open step'' was thus false; the correct ideal is the
  two-sided \lean{MRtwo}.
\item \textbf{The formal $\C[[q]]$ route says nothing about the upper half-plane.} File
  \leanpath{VOA/Modularity/FormalSpecialization.lean}:
\begin{leancode}
theorem no_algHom_powerSeries_X_ne_zero (φ : PowerSeries ℂ →ₐ[ℂ] ℂ) {q₀ : ℂ} (hq : q₀ ≠ 0) :
    φ PowerSeries.X ≠ q₀ := by
\end{leancode}
  A $\C$-algebra homomorphism $\C[[q]]\to\C$ must send $q$ to $0$ (if $q\mapsto q_0\ne0$, then
  $1-q/q_0$, a unit of $\C[[q]]$, would map to $0$). So ``evaluation at $q_0=e^{2\pi i\tau}$''
  does not exist on formal power series and no Nakayama argument over $\C[[q]]$ reaches a point of
  the upper half-plane.
\item \textbf{\lean{BlockFiniteDim.zhu_modularity_of_analytic_finrank_le} is vacuous}: its
  hypothesis \lean{hle} is refuted by \lean{not_finrank_analyticSlashAmbient_le} in
  \leanpath{VOA/Modularity/SlashAmbientObstruction.lean}. (The underlying cause was a helper
  definition returning a junk value on inhomogeneous vectors.)
\item \textbf{\lean{qInjOnBottomResidual} is false} for $\dim A(V)>1$ (free-boson
  counterexample, certificate \lean{qInjOnBottom_false_of_balancing_ne_bot} in
  \leanpath{VOA/Modules/ZhuVermaAnArchSdBalancedBottom.lean}).
\item \textbf{\lean{SimpleAModsRealized} as first typed was false}: it forced every simple
  $A(V)$-module to be finite-dimensional, false for $A(V^k(\mathfrak{sl}_2))\cong U(\mathfrak{sl}_2)$;
  now scoped under finite-dimensionality of $A(V)$.
\item \textbf{\lean{FZModeDescent} is false} (free-boson counterexample: only the zero mode
  descends to the Frenkel--Zhu quotient).
\item \textbf{Zhu's Lemma 4.4.4 in ``unitriangular'' form is false} (it would force
  $Z(L[-n]x)=0$ for $n\ge3$); Zhu carries a Virasoro-primary hypothesis for this reason.
\item \textbf{General \lean{DiagonalFHL} is not a theorem}: full invariance for an arbitrary
  bilinear form forces symmetry conditions on $V_1$; the canonical invariant form is obtained by a
  different route.
\end{itemize}

\begin{warning}
\textbf{A mirror-image mistake.} The register also contains a \emph{retracted} dead end: an earlier
note declared the structure \lean{Sl2Presentation} unsatisfiable, judging from the \emph{name} of
an ideal (\lean{casIdeal}) without reading its one-line definition; it was satisfiable, and the
note retired live work. Lesson: before declaring a predicate
unsatisfiable, read the definition, not the name; a false ``do not attempt'' is as harmful as a
false ``done''.
\end{warning}

\subsection{Conclusions satisfied by trivial objects}

\subsubsection{The square-zero TRAP}\label{ch15:sec:trap-sqzero}

The first ``definition of done'' for $\Vnat$ in \leanpath{Notes/MonsterRoadmap.md} listed five
targets: a VOA structure on $\Vnat=V_\Lambda^+\oplus(V_\Lambda^T)^+$, central charge 24,
$V_1=0$, $\dim V_2=196884$, and graded dimension $J$. While building the reconstruction interface
(\lean{reconstruct_nat}, Chapter~\ref{ch:10}), an agent noticed that all its hypotheses hold when the
twisted$\times$twisted$\to$untwisted operator $\overline W$ is set to $0$.

\begin{proposition}[The square-zero extension]\label{ch15:prop:sqzero}
Let $U$ be a vertex operator algebra and $M$ an ordinary $U$-module with integral weights. On
$U\oplus M$ define $Y(u,z)=\mathrm{diag}(Y_U(u,z),Y_M(u,z))$ for $u\in U$ and, for $t\in M$,
$Y(t,z)(u',t')=(0,\;e^{zL_{-1}}Y_M(u',-z)t)$. Then $U\chXVsq M$ is a vertex operator algebra with
the same vacuum, conformal vector and central charge as $U$, and graded dimension
$\dim U_n+\dim M_n$. It is not simple when $M\ne0$: $0\oplus M$ is a proper ideal.
\end{proposition}

\paragraph{Proof on paper.} The products are $u\cdot u'\in U$, $u\cdot t\in M$, $t\cdot u\in M$ (by
skew-symmetry) and $t\cdot t'=0$. Locality of $Y(t)$ with $Y(t')$ holds because both map $U\oplus M$
into $0\oplus M$ and kill $0\oplus M$, so their products vanish; the remaining locality relations
are module locality and the intertwining property of the skew map $W$. The other axioms reduce to
those of $U$ and $M$. The subspace $0\oplus M$ is stable under all modes and does not contain the
vacuum.

For $U=V_\Lambda^+$ and $M=(V_\Lambda^T)^+$ this object satisfies all five original targets,
because the graded dimension depends only on the underlying graded vector space. It is not $\Vnat$.
In Lean the statement is not formalised as a separate theorem; it is recorded in
\leanpath{Notes/MonsterRoadmap.md} and in \leanpath{Notes/M8GluingPlan.md} (``Status, after M8d
part 1''), and the defence against it is built into the final theorem in two ways: the vertex map
is pinned to \lean{Ynat VnatAxioms.wbarL} with the genuine adjoint $\overline W$, proved nonzero, and
the conclusion contains \lean{V.IsSimple}. File \leanpath{VOA/Moonshine/VnatGenuine.lean}:
\begin{leancode}
theorem leechWbar_ne_zero : leechWbar ≠ 0 := by
\end{leancode}

\begin{dictionary}
\lean{leechWbar} is the family $t\mapsto\overline W(t,z)$ of operators $(V_\Lambda^T)^+\to V_\Lambda^+$
defined as the adjoint of the skew intertwiner $W$ with respect to the invariant forms
(Chapter~\ref{ch:09}). The proof exhibits a nonzero coefficient: for weight-2 twisted states,
$\overline W(t)_{(3)}t'=\mathrm{TPForm}(t',t)\,\mathbf 1$, and the form \lean{TPForm} is
nondegenerate. Nonvanishing of $\overline W$ alone would not exclude every degenerate structure; it is
\lean{IsSimple} that does, via nondegeneracy of the invariant form $B^\natural=B_V\oplus\mathrm{TPForm}$
(\lean{SimpleNat.vnatVOA_isSimple}, Chapter~\ref{ch:10}), which fails for the square-zero extension
because the $\overline W$-block of invariance is exactly the adjunction $\overline W=W^\dagger$.
\end{dictionary}

\begin{physics}
The square-zero extension is a chiral algebra in which the twisted-sector states never fuse
back into the untwisted sector: $\mathrm{tw}\times\mathrm{tw}\to 0$. It has the right partition
function $J(\tau)$ but it is not a unitary CFT: its two-point functions of twisted states vanish
($\langle t|t'\rangle$ would be computed by $t\cdot t'$ landing on the vacuum), so the twisted states
are null. A physicist's check that the theory is non-degenerate is exactly the nondegeneracy of the
Zamolodchikov two-point function, which is what simplicity encodes here. The lesson is a classic
one: the partition function does not determine the operator algebra.
\end{physics}

\begin{status}
\lean{leechWbar_ne_zero} and \lean{SimpleNat.vnatVOA_isSimple} are proved; the latter was first
conditional on \lean{htt} (twisted$\times$twisted locality), which is discharged in the final
theorem: for the final structure, \lean{V.IsSimple} is the last clause of \lean{VOA.Moonshine.moonshine}
itself (\leanpath{VOA/Moonshine/TopPairs.lean}), which has no hypotheses. The square-zero proposition itself is argued on paper only.
\end{status}

\begin{warning}
\textbf{Lesson.} For every target ask what else satisfies it. Numerical invariants (dimensions,
characters) do not determine algebraic structure. ``Done'' for $\Vnat$ requires simplicity, or
equivalently a nondegenerate invariant form.
\end{warning}

\subsubsection{The $\varphi=0$ trap and the hits\_cyclic clause}\label{ch15:sec:phi0}

The triality $\sigma$ on the twisted sectors was reduced (Chapter~\ref{ch:12}) to top-level data:
a linear map $\varphi$ from the weight-2 space of $V_3^+$ to $T_s^+$ intertwining zero modes. File
\leanpath{VOA/Moonshine/SigmaExistLeech.lean}:
\begin{leancode}
structure TopSigmaV3 (s : ZMod 2) (φ : VPlus →ₗ[ℂ] TP) : Prop where
  /-- `φ` maps the top level of `V₃⁺` into `T_s⁺`. -/
  maps_frame : ∀ x ∈ OmegaV3, φ x ∈ TPframe s
  /-- `φ` maps the top level into weight 2. -/
  maps_weight : ∀ x ∈ OmegaV3, φ x ∈ plusModule.weight (2 : ℂ)
  /-- Zero-mode intertwining: `φ (o(a) x) = o(σ a) φ x` for every homogeneous `a ∈ V₀⁺`. -/
  zero_mode : ∀ (k : ℤ) (a : V0Sub.carrier), (a : VPlus) ∈ leechVOAPlus.grading k →
    ∀ x ∈ OmegaV3, φ (actV a (k - 1) x) = actT (sig a) (k - 1) (φ x)
  /-- `φ` hits a cyclic vector `h₀(−½)𝟙 ⊗ t₀` of `T_s⁺`. -/
  hits_cyclic : ∃ t₀ : TLeechT, t₀ ∈ leechFrameT s ∧ t₀ ≠ 0 ∧ ∃ x ∈ OmegaV3, φ x = x0T t₀
\end{leancode}
\begin{warning}
Without the last clause, $\varphi=0$ satisfies the first three, and the existence theorem
\lean{sigmaV3T_exists_of_top} would produce the zero intertwiner, from which nothing about $\sigma$
follows. \lean{hits_cyclic} (its image contains a nonzero cyclic vector $h_0(-\frac12)\mathbf
1\otimes t_0$) excludes this. Lesson: every existence-by-reduction statement needs a
non-degeneracy clause, and the referee should look for it.
\end{warning}

\subsubsection{A sign that is a condition, not a normalisation}

The data of $\sigma$ (\lean{SigmaData}) leave free $\Phi\mapsto c\Phi$, $\Psi\mapsto c^{-1}\Psi$ and
$\Theta\mapsto-\Theta$. One might think every scalar condition can be met by normalising. The
condition \lean{ThetaTop} says $\mathrm{TPForm}(\Theta\cdot,\Theta\cdot)=\mu\,\mathrm{TPForm}$ with
$\mu=+1$; since $\Theta^2=1$ and the form is unique up to scalar, $\mu=\pm1$, and $\mu$ is
\emph{invariant} under $\Theta\mapsto-\Theta$. If $\mu$ were $-1$, both eigenspaces of $\Theta$
would be isotropic, the trace of $\Theta$ on the top would be $0$, and the key identity would be
unsatisfiable. Numerically $\mu=1$ to $4\times10^{-12}$ and $\tr\Theta_J=-112\ne0$
(Section~\ref{ch15:sec:num-sigmakey}); in Lean $\mu=+1$ is proved (\lean{thetaTop_mkData}).
Similarly the final identity holds for $\varepsilon=+1$ (\lean{topPairs_one}); that it fails for
$\varepsilon=-1$ is argued on paper only (it would force $\Theta_0(o(x_n)\Phi_0x_n)=0$).

\subsection{Wrong formulas and conventions}

\subsubsection{The unmodified triality of Dolan--Goddard--Montague: 6372 versus 2324}\label{ch15:sec:6372}

The triality $\sigma$ on the Niemeier lattice VOA $V_N$, $N=N(A_1^{24})$, is built from the 24
copies of $\widehat{\mathfrak{sl}}_2$ at level one. Dolan, Goddard and Montague give it as
$\sigma=\prod_j\exp\bigl(\tfrac{i\pi}{\sqrt2}(J_0^{j1}+J_0^{j2})\bigr)$ (their (7.5)) and, for
$d/8$ odd, a modification inserting a sign $z_l=(-1)^{x_l}$ in one coordinate (their (7.8)). The
docstring of \leanpath{VOA/Moonshine/TrialityGolay.lean} records a trace computation that decides
between them.

\begin{proposition}[Trace certificate]\label{ch15:prop:trace}
On $(V_0^+)_2$, $\dim=49428$,
$\tr(\sigma|(V_0^+)_2)=300+8\sum_{c\ \text{octad}}(-1)^{c\cdot u}$, where $u=0$ for (7.5) and
$u=e_l$ for (7.8). Hence $\tr=6372$ for (7.5) and $\tr=2324$ for (7.8).
\end{proposition}

\paragraph{Proof on paper.} $(V_N)_2$ decomposes under $\mathfrak{sl}_2^{24}$ into Fock and root
states and the $759\cdot2^8$ octad states $e^{n}$, $n=\pm1$ on an octad. Projecting to the fixed
points of $\langle\chi,\Theta\rangle$ with $\tfrac14\sum_{k}$ (legitimate because $\sigma$ normalises
this group) and using the $SU(2)$ characters of $\sigma_j$ on spin $\frac12$ and spin $1$, the
non-octad contribution is $300$ and each octad contributes $8(-1)^{c\cdot u}$. There are 759 octads,
so $u=0$ gives $300+8\cdot759=6372$. Each point lies in 253 octads, so for $u=e_l$ the sum is
$(759-253)-253=253$, giving $300+8\cdot253=2324$.

Now suppose $\sigma$ were the restriction of an automorphism $\tilde\sigma$ of $\Vnat$, i.e.\ (by
FLM) of an element of the Monster $\mathbb M$, with $V_0^+=(\Vnat)^E$, $E\cong2^2$. Then
$\tr(\sigma|V_0^+{}_2)=\frac14(\tr\tilde\sigma+\tr\tilde\sigma k_0+2\tr\tilde\sigma k_1)$ on
$\Vnat_2=1+196883$, with $\tilde\sigma,\tilde\sigma k_0$ involutions and $\tilde\sigma k_1$ of order
4. The largest values of the Monster characters on $196883$ for such elements are $4371$ (class 2A)
and $275$ (class 4A), so the trace is at most $\frac14(4372+4372+2\cdot276)=2324$. Thus 6372 is
impossible and 2324 saturates the bound.

\begin{physics}
The trace of a symmetry on a weight space is a twined partition-function coefficient: it is what
a physicist would compute to identify which symmetry of the Moonshine CFT an operator is. Here a
character-table bound from the Monster acts as a consistency check on a formula copied from a
paper: the unmodified rotation of the 24 $SU(2)$ currents is a symmetry of the Niemeier theory but
not one that extends to the Moonshine theory.
\end{physics}

\begin{status}
Argued on paper (in the file docstring) and checked numerically: with the real Golay code the
weight-2 model (\leanpath{scripts/moonshine/numerics/triality/}) reproduces $\dim(V_0^+)_2=49428$ and
the traces 6372 and 2324, and at module level the unmodified $\sigma$ maps $V_3^+$ to the wrong
twisted summand (Section~\ref{ch15:sec:num-triality}). The Monster bound uses the Monster character
table, which is cited, not formalised. In Lean the library's \lean{golaySigma} is defined in the
(7.8) form; nothing in Lean depends on the bound.
\end{status}

\subsubsection{Sign conventions that each turned out to be necessary}

The numerical model of Section~\ref{ch15:sec:num-e8} tested the twisted$\times$twisted locality
relation (4.23) in the library's conventions and then broke one convention at a time. Each of the
following ``simplifications'' makes (4.23) fail with relative residual $1.7$--$2.0$, and each is
therefore recorded as a dead end.
\begin{enumerate}
\item \textbf{The $\varepsilon(\alpha,\alpha)$ sign on $T$.} The invariant form on the 4096-dimensional
  module $T$ must satisfy $B(e_\alpha t,t')=(-1)^{\langle\alpha,\alpha\rangle/2}B(t,e_\alpha t')$
  (File \leanpath{VOA/Moonshine/TForm.lean}, \lean{TForm_eT}), the factor $(-1)^{L_0}$ of the
  contragredient. The ``group'' sign $\varepsilon(\alpha,\alpha)$ (transpose $=$ inverse) differs from
  it by a basis-dependent character, and the corresponding form is symmetric or \emph{alternating}
  according to an Arf invariant that depends on the chosen basis: for random bases of $E_8$ the
  count $\#\{x:x\cdot Mx=0\}$ is $136$ or $120$
  (\leanpath{VOA/Moonshine/TFormLattice.lean}).
\item \textbf{The positive Fock form.} The invariant form on $S[\theta]$ must satisfy
  $h(m)^{t}=-h(-m)$ (\lean{fockForm}, File \leanpath{VOA/Moonshine/InvFormFock.lean}); the positive,
  Hermitian-looking form $h(m)^t=+h(-m)$ transposes the dressing $E^-E^+(v)$ to $E^-E^+(-v)$ and
  breaks invariance (sympy check recorded in the file).
\item \textbf{A constant charge sign.} On $V_L$ the form is
  $B(p\otimes e^\alpha,q\otimes e^\beta)=\delta_{\alpha+\beta,0}\,c(\alpha)F(p,q)$ with
  $c(\alpha)=(-1)^{\langle\alpha,\alpha\rangle/2}\varepsilon(\alpha,\alpha)$, which is \emph{not} constant
  because the library's cocycle does not satisfy $\varepsilon(\gamma,\gamma)=(-1)^{\langle\gamma,\gamma\rangle/2}$.
  File \leanpath{VOA/Moonshine/NondegLattice.lean}:
\begin{leancode}
noncomputable def chSign (α : L.carrier) : ℂ :=
  (-1 : ℂ) ^ (LatGen.halfNorm L α) * latticeCocycle L α α
\end{leancode}
  Setting $c\equiv1$ breaks (4.23).
\item \textbf{Dropping $(-1)^w$ in $\overline W$.} The adjoint formula carries a weight sign; dropping it
  breaks (4.23).
\item \textbf{The super sign.} The anticommutator version of the twisted$\times$twisted relation
  fails (residual $\approx2$): with the library's signs the extension is a VOA, not a vertex
  superalgebra.
\end{enumerate}

\begin{physics}
These are the familiar sign headaches of lattice CFT: the cocycle $\varepsilon$ that makes vertex
operators $e^{i\alpha\cdot X}$ mutually local, the BPZ conjugation $(-1)^{L_0}$ and $z^{-2L_0}$ of a
primary field, and the reality condition $a_n^\dagger=a_{-n}$ versus the bilinear transpose
$a_n^t=-a_{-n}$ for a chiral boson $\partial X$ (the bilinear form differs from the Hermitian one by
$X\mapsto -X$ on the antilinear side). Each sign controls whether the twisted-sector OPE is
associative; getting one wrong gives a theory that is not local.
\end{physics}

\subsection{Ambiguities}

\subsubsection{The golayE ambiguity}\label{ch15:sec:golayE}

The automorphism $\sigma$ of $V_0^+$ (\lean{v0Sigma}) is transported from $V_N$ through an
isomorphism that involves a character \lean{golayE} chosen by \lean{Classical.choose} from an
existence statement. The specification determines only twice the character.

\begin{theorem}[golayE is not determined]\label{ch15:thm:golayE}
There are two characters $e_1,e_2$ of $\Lambda_0$ with values in $\C^\times$, both satisfying the
specification $\psi=g_e\,\theta\,g_e^{-1}$, whose values at the preimage $w_0$ of $8e_0$ differ by a
sign.
\end{theorem}
File \leanpath{VOA/Moonshine/TrialityTwSigma.lean}:
\begin{leancode}
theorem golayE_not_determined :
    ∃ e₁ e₂ : lam0.carrier →+ Additive ℂˣ,
      (∀ x, lam0Psi x = charLat lam0 e₁ (thetaLat lam0 (charLat lam0 (-e₁) x))) ∧
      (∀ x, lam0Psi x = charLat lam0 e₂ (thetaLat lam0 (charLat lam0 (-e₂) x))) ∧
      ∃ w₀ : lam0.carrier, inclΛ w₀ = frameW ∧ charVal lam0 e₂ w₀ = -charVal lam0 e₁ w₀ := by
\end{leancode}
\begin{dictionary}
\lean{lam0.carrier →+ Additive ℂˣ} is a group homomorphism from $\Lambda_0$ to the multiplicative
group $\C^\times$ (written additively); \lean{charLat lam0 e} is the induced automorphism $g_e$ of
$V_{\Lambda_0}$ ($e^\alpha\mapsto e(\alpha)e^\alpha$); \lean{thetaLat} is $\theta$; \lean{lam0Psi} is
the target involution $\psi$; \lean{inclΛ w₀ = frameW} says $w_0$ maps to $8e_0$ in the Leech
lattice; \lean{charVal} evaluates. The proof takes $e_2=e_1+e_\beta$ with $e_\beta=(-1)^{\langle\beta,\cdot\rangle}$,
$\beta=(5,1^{23})$, which has $2e_\beta=0$ and $\langle\beta,8e_0\rangle$ odd.
\end{dictionary}
\begin{physics}
Changing the choice conjugates $\sigma$ by the symmetry $g_\beta\oplus(1\otimes e_\beta)$ of the
Moonshine CFT, which exchanges the two twisted subsectors $T_0^+$ and $T_1^+$. Physically the
question ``which twisted sector does triality map the untwisted sector $V_3^+$ to'' has no
convention-independent answer until a choice is made; both answers describe isomorphic theories.
\end{physics}
\begin{status}
Unconditional; \texttt{\#print axioms} re-run: standard axioms. Consequence: downstream statements
are phrased with ``$\exists s$'' (e.g.\ \lean{topSigma_both}); one must not try to prove that
$\sigma(V_3^+)$ is \lean{TPframe 0} (or \lean{1}) from the definitions.
\end{status}

\subsubsection{The non-computable cocycle, and why it did not matter}

The Leech cocycle is \lean{latticeCocycle} over a basis obtained from \lean{Module.Free.chooseBasis},
a choice; and the oscillator frame is the $\sqrt G$ frame, not the Golay frame. So no explicit sign
$\varepsilon(\alpha,\beta)$ is computable, and every identity must be proved symbolically. During the
top-level work this was flagged as a design issue: an explicit Golay-frame model transported via
functoriality might simplify the remaining proofs. The agent who completed the triality reported,
and the supervisor confirmed by inspection of the finished proofs, that no concrete cocycle value
was ever needed: everything reduced to bimultiplicativity, the commutator
$\varepsilon(\alpha,\beta)\varepsilon(\beta,\alpha)=(-1)^{\langle\alpha,\beta\rangle}$, and Golay-code facts.
The refactor was not done. Related design point: comparing cocycles of a sublattice and of the
ambient lattice needs a $\Z/4$-valued coboundary; a $\pm1$-valued one would commute with $\theta$
only up to a character.

\subsubsection{Two lattice types}
There are two structures called \lean{IntegralLattice}: the vendored \lean{_root_.IntegralLattice} and
\lean{VOA.IntegralLattice}; inside \lean{namespace VOA} the bare name means the latter. The bridge
\lean{IntegralLattice.ofEbeling} transports between them. A wrong resolution of a shadowed name
typechecks and yields an unprovable goal; the same trap is recorded library-wide for
\lean{mode_iterate} and \lean{mode_commutator}.

\subsection{Circular and insufficient reductions}\label{ch15:sec:circ}

\begin{enumerate}
\item \textbf{The vacuum trick for $(V_\Lambda^T)^+$ is circular.} The untwisted sectors are cyclic
  by the ``vacuum trick'' (a Borcherds identity applied to $u_{(n)}\mathbf 1$). The analogous route
  for the twisted sector needs the $\Vnat$ Borcherds identity for three twisted vectors, which is
  (4.23), the very relation being proved. Commutator formulas with one untwisted argument do not
  imply it. Twisted cyclicity (\lean{cyclic_TPframe}) was proved by an explicit route inside the
  twisted module.
\item \textbf{\lean{YnatT_eq_of_creates} is circular for the key identity.} File
  \leanpath{VOA/Moonshine/M8h0Gap.lean}:
\begin{leancode}
theorem YnatT_eq_of_creates (hspan : ModeSpan g) (x : TP) (Ψ : FormalDist (End ℂ Vnat))
    (hcov : DCovariant Tnat Ψ) (h0 : Ψ 0 vacNat = ((0 : SkewVP), x))
    (hneg : ∀ q : ℤ, q < 0 → Ψ q vacNat = 0) (hU : ∀ v : SkewVP, Local Ψ (YnatU v))
    (hT : ∀ j, Local Ψ (YnatT leechWbar (g j))) : YnatT leechWbar x = Ψ := by
\end{leancode}
  It identifies a field with $Y^\natural(0,x)$ if the field is local with \emph{all} $Y^\natural(g_j)$,
  twisted generators included (\lean{hT}). Applied to $\Psi=\sigma Y^\natural(v)\sigma^{-1}$, the
  hypothesis \lean{hT} is twisted$\times$twisted locality, i.e.\ the goal. The replacement uses
  locality with $V_0^+$ only. File \leanpath{VOA/Moonshine/M8hKeySigma.lean}:
\begin{leancode}
theorem eq_zero_of_V0local (F : FormalDist (End ℂ Vnat))
    (hloc : ∀ u : V0Sub.carrier, Local F (YnatU (u : VPlus)))
    {β : leech.carrier} (hβ : β ∉ lambda0) (t : ZMod 2 → TLeechT)
    (ht : ∀ s, t s ∈ leechFrameT s) (ht0 : ∀ s, t s ≠ 0)
    (hvac : ∀ p, F p vacNat = 0) (h3 : ∀ p, F p (frameX0P β, 0) = 0)
    (hT : ∀ s p, F p (0, x0T (t s)) = 0) : F = 0 := by
\end{leancode}
  Similarly, using $V_3^+$-locality of $L_v=\sigma Y^\natural(v)\sigma$ is circular: it is equivalent to
  the identity for $b\in V_3^+$.
\item \textbf{``Irreducible tops $\Rightarrow$ Zhu generation'' is circular.} The numerics showed that
  the weight-2 zero modes generate $\End(\mathrm{top})$ on every top (Section~\ref{ch15:sec:num-zhu}).
  This does not prove \lean{ZhuGenerated}: the set $S=\{a:\varphi o(a)=o(\sigma a)\varphi\}$ being
  everything is exactly the conclusion; irreducibility gives only a dichotomy (the graph of
  $\varphi$, or everything), and excluding the second case needs information about $A(V_0^+)$ beyond
  its image. What irreducibility does give is uniqueness of $\varphi$ up to scalar.
\item \textbf{$V_0^+$ is not strongly generated by its weight-2 space.} $\dim$ of the weight-3 Fock
  part is $576$, while $L_{-1}$ of weight 2 gives only $300$. So the classical lemma ``strong
  generators are Zhu generators'' does not apply; the 0-mode products $a_{(0)}b$ are essential
  (filtered commutator $a*b-b*a\equiv a_{(0)}b$).
\item \textbf{Weak generation does not reduce the top clause.} $V_\Lambda^+$ is generated under
  $n$-th products by the $e^\beta+e^{-\beta}$ (\lean{leechPlus_generated}), but $o(a_{(n)}b)$ on the
  top is not determined by $o(a)$, $o(b)$ on the top. Only Zhu generation (closure under $*$ modulo
  $O(V)$) reduces the zero-mode clause to generators.
\item \textbf{(4.23) is not decided by finitely many low-weight coefficients}: the top four-point
  function is $z_2^af(z_2/z_1)$ with $f$ an infinite series whose rationality is the content; no
  finite set of coefficients certifies it without a closed form or an ODE.
\item \textbf{Top-level vanishing does not imply that an intertwiner vanishes} without the
  Frenkel--Zhu bimodule $A(M)$: the top-level part of the key identity is necessary, not
  sufficient. Machine-checked replacement: top-level \emph{output} vanishing for \emph{all} states
  $v$ suffices (\lean{keySkew_of_keyTopV}); finiteness then comes from bimodule generation.
\item \textbf{Induction on $\wt v+\wt b$} via the defect commutator and skew-symmetry does not
  close (zero-mode reductions raise the state weight); skew-symmetry alone does not decide
  \lean{KeySkew} (it says $J$ is a fixed point of an involution, which allows components in both
  eigenspaces); the relation (4.22) gives no new information about the key block.
\item \textbf{Weight-2 quadratic left factors alone fail in charged degree 1}: the exact
  charged-vacuum formula \lean{quadF_lead_one} was needed (Fock spanning).
\item \textbf{Ground states cannot be reached from Fock-dressed sources} (weight count), so the
  lattice statement \lean{GroundLattice} is necessary for the filtered criterion, not merely
  convenient.
\item \textbf{\lean{Modules.Fusion.Intertwiner} is commutator-only.} File
  \leanpath{VOA/Modules/Fusion.lean}:
\begin{leancode}
structure Intertwiner (A : OrdVModule W M) (B : OrdVModule W N) (C : OrdVModule W P) where
  /-- The mode data `𝒴_s(u)x`. -/
  modes : IntertwinerModes M N P
  /-- Lower truncation: `𝒴_s(u)x = 0` for `Re s ≫ 0`. -/
  truncation : ∀ (u : M) (x : N), ∃ s₀ : ℝ, ∀ s : ℂ, s₀ ≤ s.re → modes s u x = 0
  /-- The defining intertwining commutator (membership in `IntertwinerSpace`). -/
  isIntertwiner : modes ∈ IntertwinerSpace A B C
\end{leancode}
  It has truncation and the commutator formula but no $L_{-1}$-derivative property and no
  associativity, so it is weaker than an intertwining operator of Frenkel--Huang--Lepowsky. Do not
  route uniqueness arguments through it; the Moonshine proof avoids it entirely. Its fusion
  coefficient is discussed in Section~\ref{ch15:sec:debt}.
\end{enumerate}

\begin{warning}
\textbf{Lesson for (E).} For each reduction write down what each step \emph{uses}, not only what it
proves. Circularity in this project always entered through a hypothesis that looked general
(``local with all generators'', ``a uniqueness lemma'') but specialised to the goal.
\end{warning}

\subsection{Combinatorial and numerical pitfalls}
\begin{itemize}
\item \textbf{Random sign sampling misses lattice exceptions.} In the ground-state step of Zhu
  generation the exceptional families $E_2=\{(\pm1)^{24}:\text{negative set}\in\mathcal G\}$ and
  $E_3=\{s+2s_ie_i:N(s)\Delta\{i\}\in\mathcal G\}$ are single sign classes indexed by Golay
  codewords; random sampling of signs found only $E_1=\{\pm4e_i\}$. $E_2,E_3$ were found by hand
  and confirmed by exhaustive enumeration. Use orbit structure or exhaustive enumeration.
\item \textbf{Do not exclude $\alpha=\beta$} in the ground product with $\langle\alpha,\beta\rangle\ge1$:
  $\gamma=2\alpha$ has $\gamma/2$ as its only partner, and an earlier \lean{GroundRel} with
  $\alpha\ne\beta$ was too weak.
\item \textbf{Low-order traces do not distinguish modules}: $\tr o(u)$ and $\tr o(u)o(v)$ agree on
  all six weight-2 $V_0^+$-modules; the $V_3^+/V_3^-$ distinction lives in long loops $n\to-n$.
\item \textbf{Ground-state rationality is not a unimodularity test}: $D_8$ also gives a rational
  four-point function.
\item \textbf{Left action alone} (no 0-modes) gives rank $0/24$ at weight 3 in the bimodule test.
\end{itemize}

\subsection{Tooling and process traps}
\begin{itemize}
\item \textbf{\lean{abbrev ExtraT} times out.} File \leanpath{VOA/Moonshine/ExtraspecialModule.lean}:
\begin{leancode}
noncomputable abbrev ExtraT : Type := Rep (cocMat L) (isUnit_det_cocMat L hEven hU)
\end{leancode}
  Instance search for the scalar action through the \lean{abbrev} times out; downstream code uses
  \lean{Rep (cocMat L) _} directly.
\item \textbf{The audit trap} \lean{-DmaxSynthPendingDepth=3} (Section~\ref{ch15:sec:trap}).
\item \textbf{\lean{ring} inside \lean{first}} silently ``succeeds'' (falls back to \lean{ring_nf});
  use \lean{ring1}. \textbf{\lean{ext} on a \lean{FormalDist₂}} descends too far; use
  \lean{FormalDist₂.ext}.
\item \textbf{Disk and pickles in numerics.} The first triality pipeline wrote multi-GB pickled
  intermediates (one stage alone about 1\,GB, about 2\,GB in total); they were deleted, and the
  KeySkew numerics were rebuilt to compute $\sigma$ block by block with a cache of about 68\,MB. Scripts
  first lived in a temporary session scratchpad and had to be copied to
  \leanpath{scripts/moonshine/numerics/} to survive; paths inside them may need adjusting.
\item \textbf{A naive rename.} When the library was renamed from \lean{MoonshineVOA} to \lean{VOA}, a case-insensitive substitution would have destroyed the genuine
  monstrous-moonshine module names (\lean{MoonshineHauptmodul}, \lean{MoonshineModule}, \ldots);
  only the exact token was substituted.
\item \textbf{Concurrent sessions in one tree} invalidated recorded builds
  (Section~\ref{ch15:sec:build}).
\item \textbf{Stale prose} (Section~\ref{ch15:sec:agents}, and Section~\ref{ch15:sec:debt}).
\end{itemize}

\subsection{Documentation debt and open library issues found while writing the book}\label{ch15:sec:debt}

The authors of the other chapters, reading every quoted declaration against its source, recorded
discrepancies. Each item below was re-checked against the source for this chapter.
\begin{enumerate}
\item \textbf{The structure called ``vertex algebra'' was not one} (\emph{resolved by renaming, 28 October 2026}). The
  library's transcription of Gannon's Definition 5.1.3(a), then named \lean{VertexAlgebra} (and its ungraded companion
  \lean{BareVA}), has no translation-covariance axiom, and VA1--VA5 imply neither covariance nor the Borcherds identity:
  a three-dimensional counterexample was machine-checked for Chapter~\ref{ch:L03} (Theorem~\ref{l03:thm:counter}) and is
  now in the library as \lean{TranslationCovariantCounterexample}. No theorem was wrong --- every Borcherds-type result
  is stated for \lean{VertexOperatorAlgebra}, whose axiom (VOA3) supplies covariance --- but the name and an old
  docstring were. The structures are now \lean{PreVertexAlgebra} and \lean{BarePreVA}; \lean{VertexAlgebra} is the
  genuine notion (a \lean{PreVertexAlgebra} with the field \lean{translation_covariant}), and
  \lean{VertexOperatorAlgebra.toVertexAlgebra} makes every VOA one. At the same time the VOA structure was layered:
  \lean{GradedVertexOperatorAlgebra} (the FLM/Lepowsky--Li notion) and \lean{VertexOperatorAlgebra} (that, of CFT type).
  Similarly the universal affine VOA $V^k(\mathfrak g)$, formerly \lean{simpleAffineVOA}, is now \lean{universalAffineVOA}
  (Chapter~\ref{ch:L23}). The super analogue \lean{SuperVertexAlgebra} has the same omission and keeps its name
  (Chapter~\ref{ch:L17}).
\item \textbf{\lean{VOA.fusionCoeff} is identically $0$} (argued on paper, not formalised). File
  \leanpath{VOA/Modules/Fusion.lean}:
\begin{leancode}
noncomputable def fusionCoeff (A : OrdVModule W M) (B : OrdVModule W N) (C : OrdVModule W P) : ℕ :=
  Module.finrank ℂ (IntertwinerSpace A B C)
\end{leancode}
  Here \lean{IntertwinerModes M N P} is the space of all families $\C\to\mathrm{Hom}(M\otimes N,P)$
  and \lean{IntertwinerSpace} the kernel of the commutator defect
  $a^P_{(p)}\mathcal Y_s(u)x-\mathcal Y_s(u)a^N_{(p)}x-\sum_j\binom pj\mathcal Y_{s+p-j}(a_{(j)}u)x$, with no
  truncation or grading condition. The defect relates $\mathcal Y_s$ only to $\mathcal Y_{s+\Z}$ and is
  invariant under the shift $s\mapsto s+c$. Hence the kernel is the product, over the classes
  $\rho\in\C/\Z$, of the spaces $K_\rho$ of solutions supported on $\rho$, and the shift by $c$ maps
  $K_\rho$ isomorphically onto $K_{\rho+c}$. Either every $K_\rho$ is $0$ (and the space is $0$) or
  all of the uncountably many $K_\rho$ are nonzero (and the space is infinite-dimensional). Mathlib's
  \lean{Module.finrank} is $0$ for infinite-dimensional spaces, so in both cases
  \lean{fusionCoeff} is $0$. So \lean{fusionCoeff} is not the fusion rule
  $N^P_{MN}$.
\item \textbf{\lean{SuperTwistedModule.jacobi}} is unrepaired (Section~\ref{ch15:sec:super}).
\item \textbf{\lean{IsLowestWeight} does not require $M_h\ne0$.} File
  \leanpath{VOA/Modules/ZeroModeMult.lean}:
\begin{leancode}
def IsLowestWeight (N : OrdVModule W M) (h : ℂ) : Prop :=
  ∀ α : ℂ, N.weight α ≠ ⊥ → ∃ n : ℕ, α = h + n
\end{leancode}
  It says every occurring weight is in $h+\N$; it does not say that $h$ occurs. Any $h$ below the
  true lowest weight by a natural number also qualifies. Statements using it as ``the lowest weight
  is $h$'' should add $M_h\ne\bot$ where the intended meaning needs it.
\item \textbf{\lean{SelfDual} is never linked to \lean{IsUnimodular}.} File
  \leanpath{VOA/Foundations/Lattices.lean}:
\begin{leancode}
def SelfDual : Prop := ∀ x : ℝ ⊗[ℤ] L.carrier, x ∈ L.dual ↔ ∃ v : L.carrier, L.toReal v = x
\end{leancode}
  and in the vendored library, File \leanpath{VOA/LatticesAndCodes/Lattice/Defs.lean}:
\begin{leancode}
def IsUnimodular : Prop := IsUnit Γ.det
\end{leancode}
  (\lean{VOA.Modules.StronglyRational} has an unrelated \lean{SelfDual} for VOAs.) No theorem in the
  library proves that the two lattice notions agree; results stated with one cannot be fed the other
  without a bridge.
\item \textbf{\lean{latMorph_injective} assumes equal rank.} The section variables of
  \leanpath{VOA/Moonshine/LatticeMorph.lean} include \lean{(hr : rank K = rank L)}; the
  functoriality $V_K\to V_L$ is therefore proved only for isometric embeddings between lattices of the
  same rank (sublattices of finite index), which is all the Moonshine construction needs
  ($\Lambda_0\subset\Lambda$, $\Lambda_0\subset N$), but not general functoriality.
\item \textbf{Stale docstrings} (\emph{resolved, October 2026}). Several Moonshine file headers said ``Not imported by the
  umbrella'' or ``not wired into the umbrella'' (for example \leanpath{VOA/Moonshine/TFormLattice.lean} and the
  header of \leanpath{VOA/Moonshine/TopPairs.lean}); a search of both files on 28~October~2026 finds no such header. Route notes contain superseded estimates and ``open'' claims, corrected in later
  ``Update''/``Status'' sections.
\item \textbf{\leanpath{README.md} was behind the library} (\emph{resolved, October 2026}: the current README no longer
  contains these sentences). It said ``It is not a formalization of the
  Moonshine module $V^\natural$ itself'' and, under Scope, ``$V^\natural$ is not formalized''; and it said
  ``The unrestricted minimal-model theorem remains open''. Both were false at the Moonshine commit:
  \lean{VOA.Moonshine.moonshine} constructs $\Vnat$, and File
  \leanpath{VOA/Examples/VirasoroMinimalRigidityComplete.lean} proves
\begin{leancode}
theorem minimal_graded_rigidity (p q : ℕ) (hp : 2 ≤ p) (hq : 2 ≤ q)
    (hcop : p.Coprime q) :
    IsRigidGraded (virasoroSimple (minimalCentralCharge p q)) := by
\end{leancode}
  with \texttt{\#print axioms} \lean{[propext, Classical.choice, Quot.sound]} (re-run for this
  chapter).
\item \textbf{\leanpath{Notes/MonsterRoadmap.md}, row of the twisted-module milestone,} called twisted associativity
  \lean{star_iterate_range} ``unused but proved'' (\emph{resolved, October 2026}: the row now says it is used); it is used in
  \leanpath{VOA/Moonshine/LatticeTwistedProduct.lean} and \leanpath{VOA/Moonshine/LatticeTwistedBase.lean}.
\end{enumerate}

\begin{physics}
The fusion coefficient $N^P_{MN}$ counts the independent ways two primaries $M,N$ can fuse into $P$
in the OPE; it is a small integer in a rational CFT. A ``count'' that is always zero is a warning
sign that the space being counted is not the space of chiral vertex operators (it lacks the
$L_{-1}$-covariance $\partial_z\mathcal Y=\mathcal Y(L_{-1}\cdot)$ and the conformal-weight constraint on
the exponents $s$). Similarly, a lowest-weight condition that does not demand a lowest-weight
\emph{state} would accept the empty spectrum.
\end{physics}

\subsection{Dead ends recorded in the later chapters (October 2026)}\label{ch15:sec:deadends-oct}

The chapters added or extended in October 2026 each record their own dead ends in \lean{warning} boxes, with the
counterexample or certificate. This subsection is an index to them, sorted by the six kinds (A)--(F) introduced at the
start of this section, with one or two sentences per entry. The certificates and the full discussion are in the cited
sections; where a finding is numerical or a paper argument only, the cited chapter says so, and so do we. The dead
ends of Chapters~\ref{ch:L28} and~\ref{ch:L29} are recorded in their own sections
(Sections~\ref{l28:sec:defects} and~\ref{l29:sec:deadends}); those of the chapter in preparation will be indexed with it.

\paragraph{(B) Hypotheses that no intended object satisfies.}
\begin{enumerate}
\item \textbf{The \lean{hdec} criterion forces $C_2$-cofiniteness} (Chapter~\ref{ch:L11},
  Section~\ref{l11:sec:genwords}, warning on page~\pageref{l11:warn:hdec}). The older surjectivity criteria for the jet
  presentation of Li's associated graded algebra assume
  $V=\operatorname{span}(\mathrm{gen}\cup\{\mathbf 1\})+C_2(V)$ for a finite family $\mathrm{gen}$. The theorems are
  correct, but the hypothesis forces $R_V=V/C_2(V)$ to be finite-dimensional, i.e.\ $V$ to be lisse. A certificate
  compiled for Chapter~\ref{ch:L11} (in a scratch file, standard axioms; not library code) proves
  \lean{hdec} $\Rightarrow$ \lean{C2Cofinite W}. Combined with the unconditional non-lisseness of
  $L_{-4/3}(\mathfrak{sl}_2)$ (Chapter~\ref{ch:L16}, Theorem~\ref{l16:thm:nonlisse}) this gives a machine-checked
  counterexample: no finite family satisfies \lean{hdec} there. On paper the same holds for every quasi-lisse VOA that is not
  lisse, including all the Deligne-series VOAs; with the currents as generators it already fails in weight two, where
  $(R_V)_2\supseteq V(2\theta)\oplus\C Q\neq0$ (Chapter~\ref{ch:L24}, Section~\ref{l24:sec:presentation}). It also fails for
  the lisse minimal models $M(2,2m+3)$, $m\ge2$, with the single generator $\omega$, because $[\omega]^2\neq0$ in
  $R_V\cong\C[T]/(T^{m+1})$ (Chapter~\ref{ch:L20}, Section~\ref{l20:sec:vir}). The replacement, used throughout, is
  spanning of Li's filtration by ordered generator words.
\item \textbf{Graded unitarity ``for some flag'' is nearly empty} (Chapter~\ref{ch:L19}, Warning~\ref{l19:warn:freeflag};
  paper argument). Because the predicate \lean{GradedUnitary B R F} quantifies over an arbitrary flag \lean{F}, any VOA with
  a nondegenerate Hermitian form on each weight space satisfies it for a suitable flag, Lee--Yang included. The Virasoro
  statement \lean{VirGradedUnitary} pins the flag and is unaffected. This complements the older defect, now repaired, in
  which clause (vi) was false for every VOA (Theorem~\ref{ch15:thm:bilpos}; Chapter~\ref{ch:L19},
  Warning~\ref{l19:warn:olddefect}).
\end{enumerate}

\paragraph{(C) Wrong formulas, and the limits of a single identity.}
\begin{enumerate}
\item \textbf{Feigin's homogeneity is false: $55$ versus $63$} (Chapter~\ref{ch:L26}, Section~\ref{l26:sec:open};
  Chapter~\ref{ch:L27}, Section~\ref{l27:sec:feigin}). The natural mechanism behind Feigin's Corollary~2 (\emph{SIGMA} 4 (2008) 070; Corollary~2.2 in arXiv v1) for the $A_2$
  total bound, that the fusion filtration is homogeneous for the grading by trivial factors, fails: at $N=2$ the degree-one
  filtration piece has dimension $55$, while the sum of the dimensions of its projections is $63$ (exact computation).
  The published repair (Kus--Littelmann) and an alternative route (Setia--Rani--Khandai) are recorded in those chapters.
\item \textbf{The heat equation alone does not give $E_4^2=E_8$} (Chapter~\ref{ch:L25}, Section~\ref{l25:sec:qseries}):
  at weight $8$ the relation contains $\sigma_7$ and no combination eliminates it; Mathlib's dimension formula is used
  instead.
\item \textbf{Star bounds by termwise domination or positive expansions} (Chapter~\ref{ch:L25},
  Section~\ref{l25:sec:rigidity}; Chapter~\ref{ch:L27}, Section~\ref{l27:sec:deadends}). Domination proofs of $\le$ are
  forced to be exact (proved on paper), the plethystic exponents of the MLDE solutions alternate in sign (observed), and
  the $G_2$ step R3 of the Macdonald pilot admits no small rational certificate; the pilot was stopped and the $G_2$ input
  stays named.
\end{enumerate}

\paragraph{(E) Circular or insufficient reductions.}
\begin{enumerate}
\item \textbf{The presentation route to the $A_2$ total bound is circular} (Chapter~\ref{ch:L26},
  Section~\ref{l26:sec:open}): the surjection from the associated graded of the fusion model is free, and its
  injectivity is the statement. In the same list: the naive truncated box overcounts ($72$ against $64$ at $m=2$); plain
  tensor induction fails by fusion associativity, so the mixed models are needed; separating functionals by charge,
  weight or $\theta$-degree fail from $N=3$; leading-term independence of Chari--Loktev words fails; ungraded counting in
  the evaluation model carries no information. Chapter~\ref{ch:L27}, Section~\ref{l27:sec:deadends}, adds two traps of
  the same kind: a toroidal mode splitting treated as an algebra map before passing to the associated graded, and a
  shift image treated as an ideal although it contains the constant $1$.
\item \textbf{The generator-symbol step of the Kus--Littelmann repair is insufficient for $D_4$} (Chapter~\ref{ch:L27},
  Section~\ref{l27:sec:feigin}). The p.~437 step infers that the symbols of the whole annihilator kill the fusion vector from the
  symbols of the defining relations. For $D_4$ at $m=2$ those symbols generate an ideal of codimension $29^2+1$, one too many,
  and an explicit element $\tilde a$ of the annihilator has a symbol outside it (computation, not Lean). The statement the step is
  meant to prove holds at $m=2$ (exactly) and $m=3$ (modulo two primes); the theorem is not known to fail.
\end{enumerate}

\paragraph{(D) and (A) Missing structure that a route needs.}
\begin{enumerate}
\item \textbf{Mathlib's Serre construction of $\mathfrak e_8$} (Chapter~\ref{ch:L23}, Warning~\ref{l23:warn:serre}): Mathlib
  does not prove the presented algebra finite-dimensional or simple, so the exceptional Lie algebras were built as lattice
  weight-one spaces and from checked structure constants.
\item \textbf{Zhu's change of coordinates as a VOA isomorphism} (Chapter~\ref{ch:L25}, Section~\ref{l25:sec:zhu}): the
  library's \lean{changeCoord} has no transformation law; a symbol calculus on the $C_2$ algebra avoids it.
\item \textbf{Heisenberg chains along short roots} for the nilpotency needed in Ravinder's kernel lemma
  (Chapter~\ref{ch:L26}, Section~\ref{l26:sec:kernelG}): replaced by general nilpotency lemmas.
\end{enumerate}

\paragraph{(F) Tooling: the cost of kernel-checked certificates.} The Deligne chapters rest on large integer certificates
checked by the Lean kernel (\lean{decide}), and their architecture is shaped by recorded failures (Chapter~\ref{ch:L23},
Warning~\ref{l23:warn:kernel}; Chapter~\ref{ch:L24}, Section~\ref{l24:sec:trust}; Chapter~\ref{ch:L25},
Section~\ref{l25:sec:files}; Chapter~\ref{ch:L27}, Section~\ref{l27:sec:deadends}). Single declarations over whole tables
reached 18--41\,GB or did not finish (a $133\times133$ sector comparison: 41\,GB and a time-out after 445\,s; one
\lean{decide} over the whole $D_4$ identity: 650\,s and 41\,GB); \lean{![…]} vector tables are slow at high index
(literal lists read with \lean{getD} are used); importing the VOA stack into check files costs about 6\,GB before any work
(light-import copies bridged by \lean{rfl} are used); a plain box enumeration of the $E_7$ roots did not finish in ten
minutes. The working pattern is many small declarations of a few thousand entries each, under a per-file memory budget.

\begin{physics}
Several of these dead ends have a direct physical reading. The failure of \lean{hdec} is the statement that a
non-rational chiral algebra such as $L_{-4/3}(\mathfrak{sl}_2)$, whose Higgs-branch-like associated variety has positive
dimension, cannot have its $C_2$ algebra spanned by finitely many states. And Feigin's $55\neq63$ is
a reminder that a filtration on a fusion product, the mathematical form of an operator product with prescribed orders
of vanishing at the insertion points, need not respect an auxiliary grading that is natural on the factors.
\end{physics}

% ---------------------------------------------------------------------------------------------
\section{The numerical evidence}\label{ch15:sec:numerics}

\subsection{Epistemic status}

The directory \leanpath{scripts/moonshine/numerics/} preserves the Python scripts and logs written
while choosing conventions and routes for the Moonshine construction. Its \lean{README.md} states
the principle: ``They are evidence, not proof; every claim used in Lean is proved in Lean.'' No
Lean declaration imports, cites or depends on any of these computations; the compiler never sees
them. Their role was strategic: to choose signs and normalisations before investing thousands of
lines in a proof, to detect wrong formulas from the literature, and to decide whether a reduction
was worth formalising. Two kinds of computation occur.
\begin{itemize}
\item \emph{Exact} computations, in integers, rationals (\lean{fractions.Fraction}) or modulo the
  prime $p=2^{31}-1$. A rank computed modulo $p$ is at most the rank over $\Q$ of an integer matrix,
  so \emph{full} rank mod $p$ certifies full rank over $\Q$; a deficit mod $p$ would not have been
  conclusive.
\item \emph{Floating-point} computations (NumPy/SciPy, double precision), reported as the maximal
  \emph{relative residual} $\|\mathrm{lhs}-\mathrm{rhs}\|/\|\mathrm{lhs}\|$ over the tested identities.
  Values near $10^{-15}$--$10^{-11}$ are rounding error; values of order $1$ (typically $2.0$, the
  residual of $\mathrm{lhs}=-\mathrm{rhs}$) mean the identity fails.
\end{itemize}
A test without a \emph{control} proves little: an identity that holds with every sign convention
tests nothing about conventions. Every floating-point experiment below therefore also reports
deliberately wrong variants that must fail.

\subsection{The \texorpdfstring{$E_8$}{E8} model of the twisted-by-twisted relation (4.23)}\label{ch15:sec:num-e8}

\emph{Files:} \lean{orb.py} (419 lines), \lean{sanity.py}, \lean{sanity2.py}, \lean{tt.py},
log \lean{tt_K3.log}. \emph{Question:} with the library's exact conventions, does the gluing of
$V_L^+$ and $(V_L^T)^+$ satisfy the last locality relation (4.23) of Dolan--Goddard--Montague?
\emph{Why $E_8$:} rank 8 makes $d/8$ odd as for Leech, so $\theta$ acts by $-1$ on $T$ exactly as in
the library, while $\dim T=16$, $\dim\mathrm{TP}_1=128$, $\dim\mathrm{TP}_2=2048$ are small, and the
orbifold is known to be $V_{E_8}$ again. \emph{Method:} the script reimplements
$Y^T(e^\alpha,z)=2^{-\langle\alpha,\alpha\rangle}z^{-\langle\alpha,\alpha\rangle/2}E^-E^+\otimes e_\alpha$, composite
states via $e^{\Delta_z}$, the basis cocycle, the forms with the signs of
Section~\ref{ch15:sec:deadends}, $W$ and $\overline W$ (the latter by summing over a dual basis), and
tests the twisted block of the commutator formula
\begin{plaincode}
W(t1)_m Wbar(t2)_n x - W(t2)_n Wbar(t1)_m x = sum_{j < wt t1 + wt t2} C(m,j) Y^T(Wbar(t1)_j t2)_(m+n-j) x
\end{plaincode}
for random $t_1,t_2,x\in\mathrm{TP}_1\oplus\mathrm{TP}_2$ and all $(m,n)$ with intermediate weight $\le3$
(34752 states of $V_L$). \emph{Sanity checks:} $Y^T(\omega)$ gives $L_0+d/16$, $L_{\pm1}$; Heisenberg
commutators (residual $\le10^{-16}$); lattice commutators for $\langle\alpha,\beta\rangle\in\{-1,0,1\}$
($\le3\times10^{-15}$). \emph{Result:} 64 identities, maximal relative residual $3.7\times10^{-14}$
(final log line ``max rel residual 3.734689360353892e-14 over 64 identities''). \emph{Controls}
(18 identities each), all failing with residual $1.7$--$2.0$: group sign $\varepsilon(\alpha,\alpha)$ on
$T$; $c(\alpha)=1$; dropping $(-1)^w$ in $\overline W$; the super (anticommutator) sign.

\begin{status}
Numerical evidence for a model lattice, in floating point; not about Leech, not part of any
proof. The Leech relation (4.23) is proved in Lean by a different route (triality).
\end{status}

\subsection{The theta check}

\emph{File:} \lean{theta_check.py}, log \lean{theta_check.log}. \emph{Question:} what lattice
input makes the ground-state four-point function rational? \emph{Method:} exact rational $q$-series
to $O(q^{80})$; $x=\theta_2^4/\theta_3^4$ computed by series division, without assuming Jacobi's
identities. \emph{Results} (verbatim from the log): the Leech counts
$1,0,0,0,196560,0,16773120,0,398034000$; ``Jacobi check $1-x=\theta_4^4/\theta_3^4$: True'';
$\Theta_{E_8}=\theta_3^8(1-x+x^2)$; and
\begin{equation*}
  \Theta_\Lambda=\theta_3^{24}\Bigl[(1-x+x^2)^3-\tfrac{45}{16}x^2(1-x)^2\Bigr],\qquad
  G_\Lambda(x)=\frac{(1-x+x^2)^3}{(x(1-x))^3}-\frac{45/16}{x(1-x)},
\end{equation*}
rational and crossing-symmetric; $D_8$ (determinant 4) also gives a rational $G$, so this does not
single out unimodular lattices.

\begin{physics}
$G_L(x)$ is the four-point function of the $\Z_2$ twist fields $\sigma(z)$ of $d$ chiral bosons
on $\mathbb{R}^d/L$, in the untwisted channel, as a function of the cross-ratio; $q(x)$ is the nome of the
double cover torus. Locality of the twisted OPE is rationality of $G$ in $x$: the four-point
function must be single-valued on the sphere.
\end{physics}

\subsection{Golay code and triality}\label{ch15:sec:num-triality}

\emph{Files:} \lean{golay.py}, \lean{golay2.py} (Golay code from quadratic residues mod 23;
4096 codewords, weights $\{0,8,12,16,24\}$), \lean{tri.py}, \lean{g.py}, \lean{adj.py},
\lean{cert.py}, and the pipeline \lean{triality/} (\lean{vn.py}, \lean{stage1}--\lean{stage3},
\lean{model.py}, \lean{tmod.py}, \lean{run4.py}, \lean{bfs2.py}, \lean{block.py}, \lean{bfsT.py},
\lean{trace.py}). \emph{Model:} the weight-2 space of $V_N$, all 196884 states, with the library's
basis cocycle; $\sigma=G(\chi z_0)G^{-1}$ (the library's \lean{golaySigma}) computed on the 1083
invariant blocks; $\Theta=\sigma\chi\sigma$; $(V_0^+)_2=(V_N)_2^{\langle\chi,\Theta\rangle}$. \emph{Results:}
$\dim(V_0^+)_2=49428$; $\tr\sigma=2324$ (and $6372$ for the unmodified $G\chi G^{-1}$),
reproducing Proposition~\ref{ch15:prop:trace}; $\Theta e^\gamma=c_\gamma e^{-\gamma}$ with $c_\gamma=1$ for
all 98256 norm-4 $\gamma\in\Lambda_0$. An explicit $\sigma$-twisted intertwiner $(V_3^+)_2\to(T_0^+)_2$
was built by breadth-first propagation from a joint eigenvector and checked on 15000 identities,
residual $4.2\times10^{-15}$; wrong target sectors give $2.0$; with the unmodified $\sigma$ the
consistent target is the other summand (the $N$-only $X_1$), a module-level confirmation of the
trace certificate; $T(1,+)\to T(1,+)$ residual $4.8\times10^{-12}$. Caveats recorded with the data:
only zero modes of weight-2 states on weight-2 spaces; $N$'s cocycle; the identification of the model
sectors with $T_0^+,T_3^+$ uses an unchecked normalisation step.

\subsection{The Zhu-generator test}\label{ch15:sec:num-zhu}

\emph{File:} \lean{zhu/zhu.py}, log \lean{zhu_0.log}. \emph{Question:} do the zero modes of
$(V_0^+)_2$ generate the zero-mode algebra on the weight-2 tops, or is a weight-4 element (as in
rank one, the ``$J$'' of $A(M(1)^+)$) needed? \emph{Method, exact:} on the $V_3^+$ top the zero
modes of $h_ah_b$ have joint eigenvalue $nn^{\mathsf T}$, giving projectors onto pairs
$\{e^n,e^{-n}\}$, and $o(e^\gamma+\theta e^\gamma)$ links pairs with $\langle n,n'\rangle=\pm2$; the graph on the
98304 vectors with these $226\,099\,200$ directed edges is connected, so the algebra $B$ contains
all matrix units. On the twisted top $B=M_{24}\otimes C$, and 400 random norm-4 vectors generate
$\Lambda_0$ (Hermite normal form covolume $8192=2^{13}$), so $C=\End(X)$. \emph{Result:} $B=\End(\mathrm{top})$
on both 49152-dimensional tops; $B$ acts irreducibly on, and separates, all 16 irreducible tops, total
dimension $1+3\cdot24^2+6\cdot2048^2+6\cdot49152^2=14\,520\,682\,177$. No weight-4 generator is
needed. (Status: exact computation plus standard Heisenberg-group representation theory; the
implication to Zhu generation needs the cited classification of $A(V_{\Lambda_0}^+)$ and was
\emph{not} used; Lean proves Zhu generation directly, \lean{zhuGenerated_weightTwo}.)

\subsection{Filtered-criterion ranks (zhugen)}

\emph{Files:} \lean{zhugen/fock_lowerprod.py}, \lean{charged_lowerprod.py}, \lean{ground_states.py},
\lean{e3_relations.py}, \lean{e4_check.py}, \lean{ll_special.py}, log \lean{zhugen.log}. Exact ranks
mod $2^{31}-1$ of the spans of lower products. Fock sector: rank 1 deficient by one exactly at
weight 4 (the rank-one $J$; a positive control showing sensitivity), full at weights 3 and 5--9;
ranks 2, 3, 4 full through weights 11, 9, 8. Charged sector: full rank for charges
$(2),(2,0),(2,0,0),(1,1,2),(2,2,2,0)$ up to Fock weight 7. Ground states: random placements of every
magnitude pattern up to norm 28 found only $E_1$; $E_2,E_3$ were found by hand and checked
exhaustively (Section~\ref{ch15:sec:circ}). All superseded by Lean proofs (\lean{groundLattice},
\lean{FockSpan.fockSpanned}, \lean{FockSpan.excitedSpanned}).

\subsection{SigmaKey normalisation}\label{ch15:sec:num-sigmakey}

\emph{Files:} \lean{sigmakey/} (\lean{pipeline.sh}, \lean{sk.py}, \lean{sk2.py}, logs). In the
$N(A_1^{24})$ weight-2 model: N1, $\Theta_J^2\propto I$ (residual $1.7\times10^{-11}$), invariant form
on the 2048-dimensional block built by eigenline propagation, $\Theta^{\mathsf T}B\Theta=\mu B$ with
$|\mu-1|<4\times10^{-12}$, and $\tr\Theta_J=-112$ (eigenvalue multiplicities $968$ and $1080$), which
independently excludes $\mu=-1$; N2, $\Phi$ is a similitude ($\kappa_0$ constant to
$6.7\times10^{-12}$); N3, the top-level key identity $\Theta o(v)\Phi(b)=\varepsilon o(b)\Phi(v)$ on 170
pairs with one global sign. \emph{Control caveat:} with $\Theta$ replaced by the identity, the N3
identity also holds up to a sign that varies from pair to pair ($\approx50/50$); $\Theta$ is what makes
the sign global, but the test does not discriminate strongly at weight 2.

\subsection{KeySkew at weights 3 and 4}

\emph{Files:} \lean{keyskew/} (\lean{lsig.py}, \lean{eng.py}, \lean{ops.py}, \lean{test_eng.py},
\lean{phi.py}, \lean{theta.py}, \lean{keyskew.py}). The model was rebuilt without multi-GB
intermediates. \emph{Model consistency} (\lean{test_eng.log}): zero modes agree with the old model to
$4.6\times10^{-16}$; $L_0$, $[L_{-1},u_{(m)}]$ to $<10^{-15}$; the commutator formula between $V_0^+$ and
$V_3$ on twisted states, 120 identities, $1.1\times10^{-15}$. \emph{Test} of
$\Theta Y^T(v,z)\Phi b=e^{zL_{-1}}Y^T(b,-z)\Phi v$, with $\Theta$ known only at weight 2, through
lowering letters:

\begin{center}
\begin{tabular}{llll}
\toprule
weight & identities & residual & controls \\
\midrule
2 & 12 pairs & $9.0\times10^{-12}$ & $\Theta\to$ scalar: does not fail (caveat) \\
3 & $\approx5300$ & $1.1\times10^{-11}$ & $-\Theta$: $2.0$; $\Theta\to\mathrm{Id}$: $\ge0.32$; drop $L_{-1}$: $\ge0.75$\\
4 & $\approx2000$ & $1.1\times10^{-11}$ & $-\Theta$: $2.0$; $\Theta\to\mathrm{Id}$: $\ge0.57$; drop $L_{-1}$: $\ge0.84$\\
4 (two-step) & $\approx1300$ & $1.1\times10^{-11}$ & $-\Theta$: $2.0$\\
\bottomrule
\end{tabular}
\end{center}
On one $T$-fibre of the weight-3 level (dimension 2624) the stacked test map has rank 323: a strong
necessary condition, not a proof of the vector identity.

\subsection{Bimodule generation}

\emph{Files:} \lean{bimodule/mod_excited.py}, \lean{mod_ground.py}, \lean{pair_lemma.py}. The
module analogue of the filtered criterion for $M=V_3^+$ over $V_0^+$ holds at weights 3, 4, 5 with
no deficit in any sector (exact mod $2^{31}-1$); ground sector exhaustive over the 4096 residues,
every sampled coset vector (301) has an $f=-1$ partner; pair lemma checked on 3937 odd magnitude
multisets up to norm 60. \emph{Control:} left action only gives rank 0 of 24 at weight 3 and 300 of
324 at weight 4. Superseded by \lean{FZGen.bimodGen_V3}.

\subsection{TopPairs pre-check}

Before the final identity was proved, the library's own symbolic formulas for both sides of
$\Theta(o(v)\Phi b)=o(b)\Phi v$ were evaluated on all 98304 odd norm-4 Leech vectors (built from
\lean{golayB}, as in \lean{leech.py}): 200{,}000 random pairs and 64 full rows, 0 failures (recorded
in \leanpath{Notes/M8GluingPlan.md}). The script for this run is not among
the preserved files; only its result is recorded.

\begin{physics}
These computations are what a physicist would do first: build a finite truncation of the
orbifold CFT (all states up to weight 2, 3 or 4), compute OPE coefficients numerically, and check
locality and associativity. Their limitation is the usual one of truncated numerics: they test
finitely many coefficients of identities between infinite series.
\end{physics}

\subsection{Summary}
\begin{center}
\begin{tabular}{p{3.0cm}p{3.6cm}p{2.6cm}p{4.2cm}}
\toprule
experiment & checks & precision & role / status \\
\midrule
$E_8$ (4.23) & convention consistency & $3.7\times10^{-14}$ & fixed three signs; evidence\\
theta check & lattice input of (4.23) & exact, $O(q^{80})$ & route analysis; evidence\\
Golay traces & which $\sigma$ & exact & rejected DGM (7.5); evidence\\
triality model & $\sigma$ on twisted tops & $4\times10^{-15}$ & existence plausible; superseded\\
Zhu test & generators of tops & exact & chose $G=(V_0^+)_2$; superseded\\
zhugen & filtered criterion & exact mod $p$ & found $E_1$--$E_3$; superseded\\
SigmaKey & normalisation & $4\times10^{-12}$ & $\mu=1$; superseded\\
KeySkew & weights 3, 4 & $1.1\times10^{-11}$ & de-risked; superseded\\
bimodule & module criterion & exact mod $p$ & de-risked FZ step 2; superseded\\
TopPairs & final identity & exact symbolic & 0 failures; superseded\\
\bottomrule
\end{tabular}
\end{center}
``Superseded'' means a Lean theorem now proves the corresponding statement.

% ---------------------------------------------------------------------------------------------
\section{Design decisions}\label{ch15:sec:design}

\subsection{Why the FLM route}

Five routes to the vertex algebra structure on $\Vnat$ were compared before the gluing began
(\leanpath{Notes/M8GluingPlan.md} \S1, with the sources actually read marked as such):
\begin{center}
\begin{tabular}{p{3.6cm}p{6.0cm}p{3.6cm}}
\toprule
route & new theory needed & estimate, verdict\\
\midrule
(A) FLM triality & Golay frame, $\sigma$, cyclicity, a Goddard-type uniqueness lemma, one key identity & 20--35k lines, medium risk; \textbf{chosen}\\
(B) direct (4.23) & closed form of $\overline W$ on ground states, theta transformation of $2^{24}$ coset thetas, half-integral locality & 20--35k, high risk; cross-check only\\
(C) framed VOA (Miyamoto; Dong--Griess--H\"ohn) & Ising modules and fusion rules, code VOAs, quantum Galois theory & 50k+, high risk; rejected\\
(D) simple-current extensions (Carnahan; Huang) & braided tensor categories, rationality and fusion of $V_\Lambda^+$ & 100k+; rejected\\
(E) 2020--2026 literature & none found that bypasses triality or tensor categories & --\\
\bottomrule
\end{tabular}
\end{center}
The decisive argument for (A) was reuse: every ingredient is an explicit operator on an explicit Fock
space, which is how the library had already built general lattice VOAs. The plan packaged FLM's
argument through the reconstruction theorem: generators $V_\Lambda^+\cup T_0^+$, pairs involving an
untwisted field from general theory, twisted pairs by $\sigma$-conjugation; this removes FLM's use
of the group $2^{1+24}.\mathrm{Co}_1$. In execution the key identity was reduced further: $\sigma$ on
the twisted sectors by a top-level argument plus Zhu generation (no Verma modules), the three-point
identities to one self-coupling identity, that identity by a Frenkel--Zhu bimodule argument to a
finite weight-2 identity, and that identity to explicit Golay arithmetic.

\subsection{Other decisions}
\begin{itemize}
\item \textbf{Vendoring \lean{LatticesAndCodes}} (decision of 6 October 2026): the 25-file import
  closure of \lean{Leech.Defs}, \lean{Leech.RepCount}, \lean{Leech.RepCountHigher} copied
  byte-for-byte from \lean{Fedor-Popov/LatticeAndCodes} at commit
  \lean{da62b9e2f44a04888abf6f797a429aa8ecc1d6dd}, the only edit the import prefix (verified by
  \lean{cmp}); the package depends on Mathlib alone and \leanpath{lakefile.toml} is unchanged. Leech
  uniqueness is not vendored (not in the closure).
\item \textbf{$T$ abstractly}: the 4096-dimensional module by Wedderburn--Artin from the central
  simplicity of $\C^\varepsilon[\Lambda/2\Lambda]$, not by an explicit Golay model; the splitting
  $T=T_+\oplus T_-$ and the invariant form were obtained without the model.
\item \textbf{Keeping the cocycle symbolic} (previous section).
\item \textbf{Avoiding \lean{VOA.pbw}}: the Dong--Jiang generalised Verma module carries a
  \lean{BottomProjection} supplied by the then-axiom (while \lean{VOA.pbw} was an axiom); the top-level graph argument
  (\lean{annCompat_of_top}) and Zhu's easy direction (\lean{ZhuModule} is \lean{pbw}-free) replaced it.
\item \textbf{Frenkel--Zhu instead of fusion rules}: Abe--Dong--Li fusion rules for $V_L^+$ would
  have needed the classification and tensor-product theory; the bimodule argument needed only the
  commutator formula and weight-2 data, avoiding the missing intertwining Jacobi identity for $W$.
\item \textbf{The minimal definitional repair} of \lean{TwistedModule}, and ``$\exists s$'' statements
  instead of fixing \lean{golayE}.
\end{itemize}

\subsection{The literature and online search}
\begin{itemize}
\item \textbf{H\"ohn's Lean ATLAS} (repository \lean{Moonshine-in-Kansas/atlas}) formalises
  $\mathrm{Co}_0$, $\mathrm{Co}_1$ and $M_{24}$; its licence forbids copying
  into this Apache-2.0 repository, so nothing was taken from it. It is the natural partner for the
  Monster problem below.
\item \textbf{Carnahan's vertex operators in Mathlib}: Mathlib 4.30.0 contains
  \leanpath{Mathlib/Algebra/Vertex/VertexOperator.lean} and \lean{HVertexOperator.lean} (vertex
  operators as maps to Laurent series; its TODO lists locality and residue products). The library
  predates and does not use them; it has its own formal distributions.
\item \textbf{H\"ohn--Seysen}, arXiv:2508.01037, uses $\Vnat$ and the structure of $V_\Lambda^{++}$
  (16 irreducibles, the triality as an $S_3$) but does not construct $\Vnat$.
\item Huang, arXiv:1902.09734, gives twisted$\times$untwisted operators in general, not
  twisted$\times$twisted. The FLM book was not accessible online; DGM (hep-th/9410029) \S\S3--7 was
  read in full; DGM's own \S7 derives $\sigma$ from analytic input and is circular for this purpose.
\end{itemize}

% ---------------------------------------------------------------------------------------------
\section{Reproducibility}\label{ch15:sec:repro}

\begin{itemize}
\item \textbf{Commit:} this edition documents the release commit \texttt{\bookcommit} (front matter;
  Appendix~\ref{ch:provenance}). The original verification run of the Moonshine construction was made at the Moonshine
  commit (``Construct the Moonshine module $V^\natural$ (FLM) unconditionally'', 7 October 2026, 268 files changed, 70349
  insertions, 8 deletions relative to its parent); the records below (job count, probe output) are those of that run.
\item \textbf{Toolchain:} \lean{leanprover/lean4:v4.30.0} (file \leanpath{lean-toolchain}), Mathlib
  \lean{v4.30.0} (\leanpath{lakefile.toml}); \lean{elan} at \texttt{\textasciitilde/.elan/bin}.
\item \textbf{Commands}, from the repository root:
\begin{plaincode}
export PATH="$HOME/.elan/bin:$PATH"
lake exe cache get     # Mathlib build cache, several GB
lake build             # aggregate build; recorded: 10343 jobs
python3 scripts/check_repository.py
python3 scripts/audit/run_axiom_audits.py
\end{plaincode}
\item \textbf{Moonshine axiom probe}: a file containing
\begin{plaincode}
import VOA
#print axioms VOA.Moonshine.moonshine
\end{plaincode}
  compiled with \lean{lake env lean -DmaxSynthPendingDepth=3}; expected output
  \lean{'VOA.Moonshine.moonshine' depends on axioms: [propext, Classical.choice, Quot.sound]}
  (reproduced for this chapter; about 7 seconds once the library is built).
\item \textbf{Numerics}: Python 3 with NumPy, SciPy, SymPy; each subdirectory's \lean{README.md}
  gives the run order; paths may need adjusting; the KeySkew cache is about 68\,MB and is recomputable
  in about 30 minutes; the original triality pipeline writes about 2\,GB.
\item \textbf{Never commit} \leanpath{.lake/}.
\end{itemize}

% ---------------------------------------------------------------------------------------------
\section{Open problems}\label{ch15:sec:open}

\subsection{Moonshine}
\begin{enumerate}
\item \textbf{$\Aut(\Vnat)\cong\mathbb M$.} There is no Monster in Lean. A realistic first form
  defines $\mathbb M:=\Aut(\Vnat)$ and proves facts about it: its order (H\"ohn--Seysen give a
  VOA-based computation of $|\mathbb M|$), the Griess algebra on $\Vnat_2$, the subgroup
  $2^{1+24}.\mathrm{Co}_1$. The library has the ingredients for the extraspecial group acting on $T$,
  and H\"ohn's ATLAS has $\mathrm{Co}_1$, but its licence forbids copying.
\item \textbf{McKay--Thompson series and Borcherds' theorem.} For $g\in\mathbb M$,
  $T_g(\tau)=\sum\tr(g|\Vnat_n)q^{n-1}$ is a Hauptmodul (Borcherds 1992). Needed: the Monster action,
  twined characters, the monster Lie algebra and the no-ghost theorem. The library's
  \leanpath{VOA/Modularity/Moonshine*} files contain abstract $c=24$ statements conditional on
  hypotheses never instantiated.
\item \textbf{Uniqueness of $\Vnat$.} FLM conjectured that $\Vnat$ is the unique holomorphic VOA of
  central charge 24 with $V_1=0$. Dong--Griess--Lam proved uniqueness under an additional frame
  hypothesis (cited, not checked); the general conjecture is open.
\item \textbf{A Moonshine audit probe} in \leanpath{scripts/audit/}, so that the maintained audit
  covers \lean{VOA.Moonshine}.
\item \textbf{Simplifications}: formalise the failure for $\varepsilon=-1$; make \lean{golayE} explicit;
  prove irreducibility of the four $V_0^+$-sectors; the full intertwining Jacobi identity for $W$;
  route (B) as an independent check.
\end{enumerate}

\subsection{Library-wide}
\begin{enumerate}
\item \textbf{Repair \lean{SuperTwistedModule.jacobi}} (Section~\ref{ch15:sec:super}); until then
  theorems with such hypotheses for $N\ge2$ are vacuous.
\item \textbf{Uniform negative-admissible $\mathfrak{sl}_3$ graded rigidity is open} (an open residual of the project's
  $\mathfrak{sl}_3$ work). Recorded remaining boundary: the finite count/product
  exhaustion and normalisation between total PBW root counts and real-root/descendant factors, the
  higher-rank horizontal quotient or parabolic vacuum determinant transfer, and simple vacuum
  multiplicity/nonzero tangent. A removal-bijection draft is outside the build.
\item \textbf{Graded unitarity} (Arabi Ardehali--Beem--Lemos--Rastelli): the Lean layer formalises the definition and
  reductions; the full conjecture is open (beyond the cases proved in Chapter~\ref{ch:L19}). Chapters~\ref{ch:L28}
  and~\ref{ch:L29} reduce it for the minimal models $M(2,2n+1)$ to the hypothesis \lean{hA}, graded unitarity of the
  boundary admissible $L_k(\mathfrak{sl}_2)$, which is open for every $n$.
\item \textbf{PBW in Mathlib}, which would replace the library's own proof of \lean{VOA.pbw} (formerly a project axiom, now a theorem; Section~\ref{ch15:sec:trust}).
\item \textbf{Fusion rules}: a genuine intertwining-operator space with $L_{-1}$-property and weight
  constraints, so that the fusion coefficient means $N^P_{MN}$ (Section~\ref{ch15:sec:debt}).
\item \textbf{Documentation debt} (Section~\ref{ch15:sec:debt}): \lean{IsLowestWeight}, \lean{SelfDual} versus
  \lean{IsUnimodular}. (The README, the stale Moonshine headers and the roadmap row were corrected in October 2026.)
\item \textbf{Upstreaming to Mathlib}: the level-one structure theorem $M(\mathrm{SL}_2(\Z))=\C[E_4,E_6]$
  (\lean{levelOne_coe_mem_E46}), the majorant estimate, the vendored lattice and code theory,
  the Wedderburn--Artin dimension count for extraspecial groups, and formal calculus aligned with
  Carnahan's \lean{VertexOperator}.
\end{enumerate}

\subsection{Open problems named in the later chapters}\label{ch15:sec:open-later}

Each item below is a precise statement that the library either takes as an explicit hypothesis or does not state. All
the theorems mentioned are \texttt{\#print axioms}-clean; none of those with an open hypothesis is unconditional
(Section~\ref{ch15:sec:uncond}), and a citation of the literature in a docstring does not change that.

\begin{enumerate}
\item \textbf{The named inputs of the Deligne final theorems} (Chapter~\ref{ch:L27}, Sections~\ref{l27:sec:finals},
  \ref{l27:sec:ledger} and~\ref{l27:sec:cores}). For each type $X\in\{A_2,G_2,D_4,F_4,E_6,E_7,E_8\}$ of the Deligne series,
  the final theorem \lean{X_classicallyFree_final} says that the simple affine VOA $L_k(\mathfrak g)$ at
  $k=-1-h^\vee/6$ is classically free for the explicit presentation of Chapter~\ref{ch:L24}; for every type but $A_2$ it is
  conditional on named inputs. As recorded in Chapter~\ref{ch:L27}:
  \begin{itemize}
  \item $A_2$: none. The total-dimension bound (\lean{CVUniv.fusionTotalDim}), the Demazure realizations and the star
    identity are proved, and \lean{A2_classicallyFree_final} has no hypothesis; it is machine-confirmed: the aggregate build was
    green on 2026-10-28 and \texttt{\#print axioms} reports \lean{[propext, Classical.choice, Quot.sound]} (Chapter~\ref{ch:L27}, Theorem~\ref{l27:thm:A2});
  \item $D_4$: \lean{FusionTotalDimG D4sector} ($\le28^m$) only (\lean{D4_classicallyFree_of_inputT}). The Demazure
    realizations for all $n$ are proved in Lean (\lean{D4Char.d4_demazureRealization}; aggregate build green on
    3~November~2026, \texttt{\#print axioms} standard), and the star identity is proved (\lean{D4Wp.d4SchurEqE4} with the unclamped
    Demazure $=$ Schur identity, \lean{d4DemazureEqSchur_of_positivity}, positivity from U5 on the unclamped route).
    The total-dimension bound is Kus--Littelmann's Theorem~4.2; one step of its published proof (p.~437) is insufficient as
    written for $D_4$ at $m=2$, by computation with an explicit witness, while the theorem is not known to fail. Progress on
    the bound, not part of the claim (Chapter~\ref{ch:L27}, Section~\ref{l27:sec:feigin}): it is reduced in Lean to $\mathfrak{sl}_4$ core bounds with at
    least two adjoint sites (\lean{d4_classicallyFree_of_twoAdjointCoreBounds}), and on paper the levels with two adjoint sites
    are reduced further to two one-node bounds, open in general; an earlier route through a rank bound on a presented module
    (\lean{D4_classicallyFree_of_presentedRankBoundE}) is proved for $m\le2$;
  \item $G_2$, $F_4$, $E_6$, $E_7$, $E_8$: \lean{FusionTotalDimG Xsector} ($\le(\dim\mathfrak g)^m$), the Demazure
    realizations, and the star bound (\lean{G2StarBound}, \lean{F4StarBound}, \lean{E6StarBound}, \lean{E7StarBound},
    \lean{E8StarBound}), checked numerically to orders $q^{60}$, $q^{25}$, $q^{20}$, $q^{16}$, $q^{14}$ respectively (scoping pass of
    28~October~2026, Chapter~\ref{ch:L27}, Section~\ref{l27:sec:starinput}).
  \end{itemize}
  The inputs that remain open in Lean (Chapter~\ref{ch:L27}, Section~\ref{l27:sec:cores}) are the total-dimension bound and
  the graded Demazure realizations beyond $A_2$ and $D_4$, both published theorems not yet formalized (the former with a proof step
  flagged by the project), and the star identities for $G_2$, $F_4$ and the $E$-types, which are conjectures with numerical
  evidence. For $G_2$ the library splits the identity into \lean{G2DemazureFermionic}, which follows from published results
  (Fourier--Littelmann, Chari--Moura, Ardonne--Kedem, Di Francesco--Kedem; Chapter~\ref{ch:L27},
  Section~\ref{l27:sec:starinput}), and \lean{G2FermionicKW}, a $q$-series identity that is open (checked through $q^{60}$).
\item \textbf{The Arakawa--Kawasetsu MLDE} (Chapter~\ref{ch:L16}, Section~\ref{l16:sec:mlde}). The theorem
  \lean{VOA.IsSqPrimary.exists_monicMLDEOn_all_analyticBlocks_of_etingofFinite} gives a genuine monic MLDE, with
  coefficients in $\C[E_4,E_6]$ of the forced weights and one equation for all analytic blocks, under Etingof-finiteness.
  Its two other hypotheses were open; \emph{both are now theorems}:
  \begin{itemize}
  \item \lean{AKWeightGrading W}: weight decompositions of $\C[E_4,E_6]$ and of $\C[E_4,E_6]\otimes V$, together with
    homogeneity of Zhu's relation module for the total weight. \emph{Now proved for every VOA of the library} (CFT type):
    \lean{VOA.akWeightGrading} in \leanpath{VOA/Modularity/AKGradingTotal.lean} (faithful, by the queue-2 release audit,
    \texttt{Notes/Textbook/referee/RELEASE\_AUDIT\_QUEUE2.md});
  \item \lean{E46SerreClosed}: the Serre derivative maps weight-$k$ elements of $\C[E_4,E_6]$ to weight $k+2$.
    \emph{Now proved}: \lean{VOA.e46SerreClosed} in \leanpath{VOA/Modularity/EisensteinSerreClosure.lean} (rated ``as
    stated'' by the follow-ups audit). It is no longer an open hypothesis.
  \end{itemize}
  The combined statement is \lean{VOA.IsSqPrimary.exists_monicMLDEOn_all_analyticBlocks_of_etingofFinite_uncond},
  conditional on \lean{EtingofFinite} only (despite its suffix); its instance
  \lean{VOA.exists_monicMLDEOn_all_analyticBlocks_simpleSl2QL} is unconditional. Both are \texttt{\#print axioms}-clean
  (recorded in \texttt{Notes/release/WS-RELEASE2.md}); the aggregate build at the release commit is pending. What
  remains open here: the order is existential and unbounded (Arakawa--Kawasetsu's order~2 at $L_{-4/3}(\mathfrak{sl}_2)$
  is not proved); the blocks are Zhu's abstract one-point blocks, and ``the vacuum character satisfies a monic MLDE'' is
  a short composition of library lemmas (\lean{torusOnePoint_isAnalyticBlock},
  \lean{onePointSubexp_adjoint_of_etingofFinite}) not packaged as one theorem. Also open: quasi-lisseness itself (associated varieties, symplectic leaves, the Etingof--Schedler theorem) is not
  formalized; and the Zhu-modularity chain under Etingof-finiteness (Section~\ref{l16:sec:zhu}) has its rank hypothesis
  discharged only for the trivial VOA.
\item \textbf{\lean{AdmissiblyRational} versus Dong--Li--Mason rationality} (Chapter~\ref{ch:L22},
  Remark~\ref{l22:rem:rational}). The library's rationality asks that every weak submodule of an admissible module have
  a weak complement; Dong--Li--Mason ask for complete reducibility of admissible modules into simple admissible modules.
  \emph{One direction is now proved}, for $C_2$-cofinite $V$:
  \lean{VOA.Classification.COne.DLMRat.admissiblyRational_of_dlmRational} (\lean{hC2 : C2Cofinite W}) derives
  \lean{AdmissiblyRational W} from \lean{DLMRational W}, a faithful statement of Dong--Li--Mason rationality (every
  admissible module is an internal direct sum of graded-minimal submodules); \texttt{\#print axioms}-clean, recorded in
  \texttt{Notes/release/WS-RELEASE2.md}; aggregate build at the release commit pending. Since $C_2$-cofiniteness is part
  of ``nice'', a VOA that is nice with DLM rationality in place of \lean{AdmissiblyRational} satisfies the Lean hypotheses
  of the Gannon--Rayhaun classification. \emph{The converse} (\lean{AdmissiblyRational} $\Rightarrow$ DLM-rational) is
  open; the release auditor sketched, without checking it, a graded-projection argument (take the degree-zero component
  of the projection onto a graded submodule along a weak complement; it commutes with every mode). (The library does have
  \lean{AdmissiblyRational.exists_iSupIndep_irreducible_ordinary}, in
  \leanpath{VOA/Classification/COne/RationalAdmissibleDecomposition.lean}: a $C_2$-cofinite analogue of Dong--Li--Mason's
  Theorem~8.1(c), decomposing ordinary modules into irreducibles. It is not either implication of the comparison.)
\item \textbf{The congruence property: the Frobenius--Schur half} (Chapter~\ref{ch:L21},
  Section~\ref{l21:sec:congruence}). The Dong--Lin--Ng congruence property of a nice VOA now has named theorems with no
  hypotheses: \lean{dlnCongruenceInput_holds}, \lean{zhuCongruenceInput_holds} and
  \lean{zhu_kernelContainsGamma_at_level} (namespace \lean{VOA.Huang.CongruenceNamed},
  \leanpath{VOA/Huang/CongruenceNamed.lean}), each quantifying over \lean{Nice} VOAs. At the minimal level the kernel statement is the full
  kernel claim of Dong--Lin--Ng; ``kernel half'' below refers only to the Ng--Schauenburg input. Of Ng--Schauenburg's
  Theorem~6.8 only the kernel half is proved (\lean{nsProjectiveKernelAll_holds}, Theorem~\ref{l21:thm:ns}); the statement
  that the level is exactly the Frobenius--Schur exponent, \lean{ProjectiveCongruenceTheorem} and its universal form
  \lean{NSProjectiveCongruenceAll}, occurs only as a hypothesis.
\item \textbf{Pseudo-unitarity in Huang's rigidity and the Verlinde formula} (Chapter~\ref{ch:L21},
  Theorem~\ref{l21:thm:rigid} and Section~\ref{l21:sec:moddata}). Two-sided rigidity
  (\lean{VOA.Huang.HuangRigid.rigid_ordinary}, built from \lean{rightRigid_ordinary} and the braiding), the Verlinde
  formula, the modular data and the congruence inputs are proved for \lean{Nice} VOAs, which are pseudo-unitary. Huang's
  theorems assume only simple, CFT type, self-dual, $C_2$-cofinite and rational. Pseudo-unitarity is used in two places in
  the proofs: to obtain $V'\cong V$, and in a positivity step of the Cardy argument. \emph{The first use is now removed}:
  \lean{VOA.contra_vacuum_equiv_of_stronglyRational} gives $V'\cong V$ from self-duality, and the \lean{_of_stronglyRational}
  copies in \leanpath{VOA/Huang/VerlindeVacuumFusionSR.lean} give the vacuum-fusion chain ($N_{ab}^0=\delta_{b,a'}$ and its
  bounds) for every strongly rational VOA (Chapter~\ref{ch:L21}, Theorem~\ref{l21:thm:vacSR}; \texttt{\#print axioms}-clean,
  recorded in \texttt{Notes/release/WS-RELEASE2.md}; aggregate build pending). The second,
  \lean{CardyPositivity.vacColumn_mem_K}, keeps \lean{Nice}, and pseudo-unitarity is essential there (Lee--Yang $M(2,5)$).
  Removing pseudo-unitarity from the headline theorems therefore needs a different route through the Cardy step, and
  remains open. Related missing statement: a
  single packaged theorem that the category of ordinary modules is a modular tensor category.
\item \textbf{Smaller items from the audits.} For the Arakawa--Dai VOAs the paper's Theorem~1.1 is now one Lean theorem,
  \lean{arakawaDai_theorem_1_1}, and the naturality of \lean{actualC2ArcEquiv} is \lean{actualC2ArcEquiv_natural}; what
  remains is that $J_\infty R_V$ is realized as \lean{PresentedJet} with a universal arc, since the library has no
  $J_\infty$ functor on abstract algebras (Chapter~\ref{ch:L20}); Deligne's theorem is formalized only in
  its positive-dimension (bosonic) form (Chapter~\ref{ch:L21}).
\end{enumerate}

\begin{physics}
The open problems above are, in physical terms: whether the Deligne-series vertex algebras (several of which, for
instance the $D_4$ and $E$-type members, are the chiral algebras of four-dimensional $\mathcal N=2$ superconformal
theories; for $G_2$ no such theory is known) are classically free, i.e.\ whether the arc space of
$R_V=\operatorname{Sym}\mathfrak g/(\text{explicit quadrics})$ maps isomorphically onto $\operatorname{gr}V$ (proved for
$A_2$; conditional on named inputs for the other types). $R_V$ is not the coordinate ring of the minimal nilpotent orbit
closure: it is non-reduced, and only its reduced variety is expected to be that closure; and whether the physicist's statements ``the
characters satisfy a modular differential equation'' and ``a rational CFT has a congruence modular representation'' can be
stated and proved without the extra positivity (pseudo-unitarity) that the present proofs use. The first is a
prediction from four-dimensional physics; the library states what remains open of each as named, checkable hypotheses.
\end{physics}

% ---------------------------------------------------------------------------------------------
\section{Referee's checklist}\label{ch15:sec:checklist}

\begin{longtable}{p{4.2cm}p{4.6cm}p{2.6cm}p{2.0cm}}
\toprule
declaration & meaning & status & fidelity\\
\midrule
\lean{VOA.Moonshine.moonshine} & simple VOA on $\Vnat$, FLM block map, $c=24$, $V_1=0$, $\dim V_2=196884$, graded dim $qJ$ & unconditional, standard axioms & \S\ref{ch15:sec:record}, Ch.~\ref{ch:14}\\
\lean{VOA.pbw} & PBW: affine vacuum $\ne0$; induced bottom injective & theorem, unconditional (formerly a project axiom) & \S\ref{ch15:sec:trust}\\
\lean{zhuModularity_pinned_of_C2Cofinite} & Zhu modularity given a bottom projection & conditional on \lean{Pr} & \S\ref{ch15:sec:trust}\\
\lean{zhuModularity_pinned_of_C2Cofinite_voaPBW} & Zhu modularity & conditional on $C_2$, semisimplicity, family data; via theorem \lean{VOA.pbw} & \S\ref{ch15:sec:trust}\\
\lean{zhuBijective} / \lean{_of_voaPBW} & Zhu's map, injective (the codomain is image-defined, so only injectivity has content) & conditional on \lean{Pr} / unconditional & \S\ref{ch15:sec:trust}\\
\lean{zhuCorrespondence_of_C2Cofinite} & Zhu's bijection $\Phi W\simeq$ simple $A(W)$-modules & $C_2$-cofinite & \S\ref{ch15:sec:releaseaudits}\\
\lean{zhuModularity_canonical_of_C2Cofinite_admissiblyRational} / \lean{_of_stronglyRational} & Zhu modularity, general form & $C_2$-cofinite and admissibly rational / strongly rational; nontrivial instance \lean{zhuModularity_V2m} & \S\ref{ch15:sec:releaseaudits}\\
\lean{isSemisimpleRing_A_of_admissiblyRational} & admissibly rational $\Rightarrow$ $A(W)$ semisimple & hypothesis \lean{AdmissiblyRational} only & \S\ref{ch15:sec:trust}\\
\lean{moonshine_of_frameLocality} & $\Vnat$ given (4.23) on four pairs & conditional (historic) & \S\ref{ch15:sec:uncond}\\
\lean{jacobiAllIndices_eq_zero} & old twisted Jacobi forces $M=0$ & unconditional & Thm~\ref{ch15:thm:jacobiAll}\\
\lean{SuperTwistedModule.jacobi} & super twisted Jacobi & defective, open & \S\ref{ch15:sec:super}\\
\lean{eq_bot_of_bilinear_pos} & bilinear positivity forces $S=0$ & unconditional & Thm~\ref{ch15:thm:bilpos}\\
\lean{not_nonempty_genFields_MRbot_virasoro} & no fields over \lean{MRbot} & unconditional & \S\ref{ch15:sec:deadends}\\
\lean{no_algHom_powerSeries_X_ne_zero} & no evaluation at $q_0\ne0$ & unconditional & \S\ref{ch15:sec:deadends}\\
\lean{leechWbar_ne_zero} & $\overline W\neq0$ (anti-TRAP) & unconditional & \S\ref{ch15:sec:trap-sqzero}\\
\lean{TopSigmaV3} & top-level $\sigma$ data with non-degeneracy & definition & \S\ref{ch15:sec:phi0}\\
\lean{golayE_not_determined} & character ambiguity & unconditional & Thm~\ref{ch15:thm:golayE}\\
\lean{YnatT_eq_of_creates} & Goddard identification (circular for key identity) & unconditional & \S\ref{ch15:sec:circ}\\
\lean{eq_zero_of_V0local} & uniqueness from $V_0^+$-locality & unconditional & \S\ref{ch15:sec:circ}\\
\lean{fusionCoeff} & ``fusion coefficient'' & always $0$ (paper) & \S\ref{ch15:sec:debt}\\
\lean{IsLowestWeight} & weights in $h+\N$ ($M_h$ may be $0$) & definition & \S\ref{ch15:sec:debt}\\
\lean{VirModule.minimal_graded_rigidity} & every minimal model is graded rigid & unconditional & \S\ref{ch15:sec:debt}\\
\lean{strictEtingofFinite_simpleSl2QL} & $L_{-4/3}(\mathfrak{sl}_2)$ not $C_2$-cofinite, finite $HP_0$ & unconditional, standard axioms & \S\ref{ch15:sec:releaseaudits}\\
\lean{IsSqPrimary.exists_monicMLDEOn_all_analyticBlocks_of_etingofFinite_uncond} & AK MLDE for Etingof-finite VOAs & conditional on \lean{EtingofFinite} only (\lean{akWeightGrading}, \lean{e46SerreClosed} proved); order unbounded & \S\ref{ch15:sec:open-later}\\
\lean{admissiblyRational_of_dlmRational} & DLM-rational $\Rightarrow$ \lean{AdmissiblyRational} & \lean{C2Cofinite}; converse open & \S\ref{ch15:sec:open-later}\\
\lean{boundary_classically_free} & Arakawa--Dai Thm 1.1 on the presented arc algebra & hypothesis $1\le m$ only & \S\ref{ch15:sec:releaseaudits}\\
\lean{rightRigid_ordinary} / \lean{rigid_ordinary} & Huang rigidity, right duals / two-sided \lean{RigidCategory} & for \lean{Nice} (pseudo-unitary) VOAs & \S\ref{ch15:sec:open-later}\\
\lean{contra_vacuum_equiv_of_stronglyRational} and \lean{_of_stronglyRational} copies & $V'\cong V$; vacuum-fusion chain & \lean{StronglyRational}; no pseudo-unitarity & \S\ref{ch15:sec:open-later}\\
\lean{dlnCongruenceInput_holds} & Dong--Lin--Ng congruence & for \lean{Nice} VOAs; FS half open & \S\ref{ch15:sec:open-later}\\
\lean{nsProjectiveKernelAll_holds} & kernel half of Ng--Schauenburg 6.8 & unconditional; FS half open & \S\ref{ch15:sec:open-later}\\
\lean{gannonRayhaun_classification} & nice $c=1$ VOAs, both directions & no hypotheses; rational $=$ \lean{AdmissiblyRational}; \lean{Type 0} & \S\ref{ch15:sec:releaseaudits}\\
\lean{X_classicallyFree_final} & Deligne series: classical freeness & $A_2$: no hypothesis (aggregate build green 2026-10-28; standard axioms); other types conditional on named inputs (Ch.~\ref{ch:L27}) & \S\ref{ch15:sec:open-later}\\
\bottomrule
\end{longtable}

% ---------------------------------------------------------------------------------------------
\section{Notes and references}\label{ch15:sec:notes}

The verification protocol is the project's own (its working rules, \leanpath{README.md}
``Axioms and trust base'' and ``Auditing a headline result'', \leanpath{Notes/MonsterRoadmap.md}
``Operating rules''). The dead ends are collected from \leanpath{README.md} ``Known dead ends'',
the project's internal session records, \leanpath{Notes/M8GluingPlan.md} (all ``Status'' sections) and the route notes
\leanpath{Notes/M8hRoute.md}, \leanpath{Notes/TrialityTwisted.md}, \leanpath{Notes/SigmaExistence.md},
\leanpath{Notes/ZhuGeneratorTest.md}, \leanpath{Notes/ZhuGeneration.md}, \leanpath{Notes/SigmaKey.md},
\leanpath{Notes/KeySkew.md}, \leanpath{Notes/BimoduleGenTest.md}.

\begin{itemize}
\item I.~Frenkel, J.~Lepowsky, A.~Meurman, \emph{Vertex Operator Algebras and the Monster}, Pure and
  Applied Mathematics 134, Academic Press, 1988 (chapters 3--4, 9--13). The book was not accessible to the project:
  the chapter, section and equation pointers to it throughout Chapters~\ref{ch:02}--\ref{ch:14} were not checked against
  the book itself and should be read as pointers taken from secondary sources (Dolan--Goddard--Montague, Gannon,
  Lepowsky--Li).
\item L.~Dolan, P.~Goddard, P.~Montague, \emph{Conformal field theories, representations and lattice
  constructions}, Comm.\ Math.\ Phys.\ 179 (1996) 61--120, arXiv:hep-th/9410029; and \emph{Conformal
  field theory of twisted vertex operators}, Nucl.\ Phys.\ B 338 (1990) 529--601.
\item C.~Dong, H.~Li, G.~Mason, \emph{Modular-invariance of trace functions in orbifold theory and
  generalized moonshine}, Comm.\ Math.\ Phys.\ 214 (2000) 1--56; \emph{Simple currents and extensions
  of vertex operator algebras}, Comm.\ Math.\ Phys.\ 180 (1996) 671--707, arXiv:q-alg/9504008.
\item I.~Frenkel, Y.~Zhu, \emph{Vertex operator algebras associated to representations of affine and
  Virasoro algebras}, Duke Math.\ J.\ 66 (1992) 123--168.
\item Y.~Zhu, \emph{Modular invariance of characters of vertex operator algebras}, J.\ Amer.\ Math.\
  Soc.\ 9 (1996) 237--302.
\item T.~Abe, C.~Dong, H.~Li, \emph{Fusion rules for the vertex operator algebras $M(1)^+$ and
  $V_L^+$}, Comm.\ Math.\ Phys.\ 253 (2005) 171--219, arXiv:math/0310425.
\item M.~Miyamoto, \emph{A new construction of the moonshine vertex operator algebra over the real
  number field}, Ann.\ of Math.\ 159 (2004) 535--596, arXiv:q-alg/9701012.
\item C.~Dong, R.~Griess, G.~H\"ohn, \emph{Framed vertex operator algebras, codes and the moonshine
  module}, Comm.\ Math.\ Phys.\ 193 (1998) 407--448, arXiv:q-alg/9707008.
\item C.~Dong, R.~Griess, C.~H.~Lam, \emph{Uniqueness results for the moonshine vertex operator
  algebra}, Amer.\ J.\ Math.\ 129 (2007) 583--609.
\item S.~Carnahan, \emph{Building vertex algebras from parts}, Comm.\ Math.\ Phys.\ 373 (2020) 1--43,
  arXiv:1408.5215.
\item Y.-Z.~Huang, \emph{A nonmeromorphic extension of the moonshine module vertex operator algebra},
  Contemp.\ Math.\ 193 (1996), arXiv:hep-th/9406190; \emph{Twist vertex operators for twisted
  modules}, arXiv:1902.09734.
\item G.~H\"ohn, M.~Seysen, \emph{The order of the Monster finite simple group}, arXiv:2508.01037 (title checked on
  arxiv.org, 28~October~2026; $V_\Lambda^{++}$ and its modules are in its \S2.2).
\item R.~Borcherds, \emph{Monstrous moonshine and monstrous Lie superalgebras}, Invent.\ Math.\ 109
  (1992) 405--444.
\item J.~van Ekeren, S.~M\"oller, N.~Scheithauer, \emph{Construction and classification of
  holomorphic vertex operator algebras}, J.\ Reine Angew.\ Math.\ 759 (2020) 61--99, arXiv:1507.08142.
\item T.~Gannon, \emph{Moonshine beyond the Monster}, Cambridge University Press, 2006.
\end{itemize}
